% Options for packages loaded elsewhere
\PassOptionsToPackage{unicode}{hyperref}
\PassOptionsToPackage{hyphens}{url}
\PassOptionsToPackage{dvipsnames,svgnames,x11names}{xcolor}
\documentclass[
  11pt,
]{article}
\usepackage{xcolor}
\usepackage[margin=2.4cm]{geometry}
\usepackage{amsmath,amssymb}
\setcounter{secnumdepth}{-\maxdimen} % remove section numbering
\usepackage{iftex}
\ifPDFTeX
  \usepackage[T1]{fontenc}
  \usepackage[utf8]{inputenc}
  \usepackage{textcomp} % provide euro and other symbols
\else % if luatex or xetex
  \usepackage{unicode-math} % this also loads fontspec
  \defaultfontfeatures{Scale=MatchLowercase}
  \defaultfontfeatures[\rmfamily]{Ligatures=TeX,Scale=1}
\fi
\usepackage{lmodern}
\ifPDFTeX\else
  % xetex/luatex font selection
\fi
% Use upquote if available, for straight quotes in verbatim environments
\IfFileExists{upquote.sty}{\usepackage{upquote}}{}
\IfFileExists{microtype.sty}{% use microtype if available
  \usepackage[]{microtype}
  \UseMicrotypeSet[protrusion]{basicmath} % disable protrusion for tt fonts
}{}
\usepackage{setspace}
\makeatletter
\@ifundefined{KOMAClassName}{% if non-KOMA class
  \IfFileExists{parskip.sty}{%
    \usepackage{parskip}
  }{% else
    \setlength{\parindent}{0pt}
    \setlength{\parskip}{6pt plus 2pt minus 1pt}}
}{% if KOMA class
  \KOMAoptions{parskip=half}}
\makeatother
\setlength{\emergencystretch}{3em} % prevent overfull lines
\providecommand{\tightlist}{%
  \setlength{\itemsep}{0pt}\setlength{\parskip}{0pt}}
% v3 header: single column, tight reference list, clickable links.
% Loaded via pandoc -H. Fixes for the six issues Havid found in the v2 PDF are
% mostly in the markdown generator; this file only carries LaTeX-side settings.
\usepackage{graphicx}
\usepackage{float}
\usepackage{caption}

% Notation typesetting. Havid asked for consistency of format (2026-09-08),
% so units and species are expressed with the packages the field uses rather
% than a hand-rolled $^2$ regex table.
%
% VERIFIED before adoption: both compile under tectonic, which auto-downloads
% mhchem.sty and chemgreek.sty from its bundle -- no manual TeX Live install.
% Verified at PDF span level too: \qty{1.3e-3}{\centi\meter\squared\per\volt
% \per\second} raises every exponent AND every minus, and \ce{[PbI4]^2-} puts
% the subscript 4 BELOW the baseline with the charge ABOVE it in one token.
% A regex table cannot script one token in two directions.
%
% mhchem version=4 pinned: version 3 changed the charge syntax, so an
% unpinned load would silently alter every ion in the paper.
\usepackage{siunitx}
\usepackage[version=4]{mhchem}

% Journal-style units: upright glyphs, thin space between value and unit,
% and a real minus in exponents. detect-weight keeps units upright inside a
% bold heading instead of inheriting the surrounding bold.
\sisetup{
  detect-all,
  detect-weight        = true,
  range-units          = single,
  % reciprocal, NOT symbol. With per-mode=symbol the electron mobility
  % rendered as "cm2/(V s)", which is typographically correct but is not the
  % convention this literature uses; every perovskite paper writes
  % "cm2 V-1 s-1" with reciprocal powers. Verified at span level: reciprocal
  % raises both the 2 and each minus sign.
  per-mode             = reciprocal,
  exponent-product     = \times,
  output-decimal-marker = {.},
  group-digits         = none
}

% Figure captions: LaTeX owns the "Figure N." label. The v2 PDF rendered
% "Figure 4: Figure 4." because the caption text ALSO began with "Figure 4.",
% written into the markdown by the generator. v3 passes a bare caption and lets
% LaTeX number it, so the label appears exactly once.
\captionsetup[figure]{labelfont=bf, labelsep=period, font=small,
                      justification=justified, singlelinecheck=false,
                      skip=6pt}

% Tighten the numbered reference list produced by pandoc's default
% CSLReferences environment.
\usepackage{etoolbox}
\AtBeginEnvironment{CSLReferences}{%
  \small
  \setlength{\parskip}{0pt}%
  \setlength{\itemsep}{0pt}%
  \setlength{\topsep}{2pt}%
}

% No section numbering from LaTeX: the generator already writes "1. Introduction",
% "2. The Certified Frontier ...", so LaTeX numbering would double it.
\setcounter{secnumdepth}{0}
\usepackage{bookmark}
\IfFileExists{xurl.sty}{\usepackage{xurl}}{} % add URL line breaks if available
\urlstyle{same}
% fallback for those not using the hyperref driver hyperxmp:
\makeatletter
\@ifundefined{xmpquote}{\newcommand{\xmpquote}[1]{#1}}{}
\makeatother
\hypersetup{
  colorlinks=true,
  linkcolor={[HTML]{1A4E8A}},
  filecolor={Maroon},
  citecolor={Blue},
  urlcolor={[HTML]{1A4E8A}},
  pdfcreator={LaTeX via pandoc}}

\author{}
\date{}

\begin{document}

\setstretch{1.05}
\begin{center}
{\LARGE\bfseries Perovskite Photovoltaics in January-June 2026: Defect passivation, interface and contact engineering\par}
\vspace{1.1em}
{\large Havid Aqoma\textsuperscript{1,2,3,4,*}\par}
\vspace{0.7em}
\begin{minipage}{\textwidth}\centering\small
\textsuperscript{1}School of Energy and Chemical Engineering, Xiamen University Malaysia, Selangor Darul Ehsan 43900, Malaysia\\[0.25em]
\textsuperscript{2}Kelip-kelip! Center of Excellence for Light Enabling Technologies, School of Energy and Chemical Engineering, Xiamen University Malaysia, Selangor Darul Ehsan 43900, Malaysia\\[0.25em]
\textsuperscript{3}Centre of Excellence for Green and Advanced Technologies in Efficient Separations (GATES), School of Energy and Chemical Engineering, Xiamen University Malaysia, Selangor Darul Ehsan 43900, Malaysia\\[0.25em]
\textsuperscript{4}Gulei Innovation Institute, Xiamen University, Zhangzhou 363200, China\\[0.55em]
*Corresponding author. Email: havidaqoma.khoiruddin@xmu.edu.my\\[2.0em]
September 11, 2026
\end{minipage}
\end{center}
\vspace{0.8em}

\subsection{Abstract}\label{abstract}

Evaluation of 3328 publications across 6 months, with 884 examined in
depth, reveals that reported efficiency records mask severe mechanistic
contradictions and pervasive validation deficits across the discipline.
Independent verification remained scarce: only 114 investigations
reported an accredited power conversion efficiency (PCE), with the
certified frontier reaching 34.85\% in monolithic tandems and 27.6\% in
single-junction cells. Rather than advancing across the observation
window, certified performance remained confined within a narrow
efficiency corridor where contradictory physical models yielded
indistinguishable macroscopic photovoltaic parameters. These certified
benchmarks were attained through competing interfacial and bulk defect
mitigation strategies. At top interfaces, reports diverged sharply on
whether non-radiative recombination suppression requires reconstructing
low-dimensional capping phases or whether such layers impose deleterious
carrier extraction bottlenecks. At buried contacts, self-assembled
monolayer (SAM) engineering alternated between static electrostatic
dipole alignment using polar co-adsorbents and mitigating dynamic
interfacial chemical redox reactions, while tandem subcell coupling
contested whether current matching demands laterally conductive
transparent conducting oxides or direct tin oxide recombination layers
across textured interfaces. Concurrently, formamidinium lead triiodide
phase stabilization conflicted between non-destructive contact-triggered
templating without interdiffusion and active cation exchange driven by
Ostwald ripening. Device scaling imposed an uncompromising area penalty,
with the maximum certified aperture reaching
\qty{655.2}{\centi\meter}\(^2\) at an efficiency of 26.8\%. This scaling
loss stems from conflicting precursor-solvent coordination kinetics
during blade and slot-die coating, where suppressing intermediate
solvent adducts conflicts with Lewis-base coordination required to
retard premature nucleation, compounded by drying non-uniformities,
localized sheet resistance penalties, and pinhole multiplication.
Operational durability evidence is severely deficient across the field:
only 39 publications reported stability lifetimes to the defined
retention threshold, with the longest reaching \qty{4100}{\hour}.
Evidence for long-term operational resilience remains thin, as operando
structural characterization demonstrates that bulky passivators designed
to relieve lattice strain induce localized lattice expansion and
mechanochemical fatigue under thermal stress, while buried SAM contacts
degrade through unresolved competition between field-driven lateral
aggregation into isolated clusters and out-of-plane molecular
desorption. Reporting practices reinforce this empirical deficit: only
39 studies implemented standardized International Summit on Organic
Photovoltaic Stability (ISOS) protocols. Abstract-level text extraction
confirms that while 52.1\% of publications stated an efficiency, only
8.5\% reported independent certification, showing that most claims rely
on uncertified dynamic sweeps without aperture masking verification,
spectral mismatch calibration, or steady-state maximum power point
tracking (MPPT). Mechanistic synthesis indicates that apparent
efficiency advances frequently reflect transient electrostatic screening
and interfacial dipole realignment rather than permanent electronic
defect remediation. Because diametrically opposed physical models
produce indistinguishable macroscopic parameters, standard
current-voltage characterization cannot separate interfacial
recombination velocity from bulk trap states. Resolving whether
operational degradation proceeds through interfacial redox, lateral
contact dewetting, or mobile point defect accumulation requires
operando, spatially resolved transient spectroscopy paired with
depth-resolved photoemission under simultaneous electrical bias,
continuous illumination, and elevated temperature, alongside
standardized ISOS aging on certified apertures.

\textbf{Keywords:} perovskite photovoltaics; buried interface; defect
passivation; wide-bandgap perovskite; tandem solar cells; operational
stability

\subsection{1. Introduction}\label{introduction}

Recent investigations into interfacial decay in metal halide perovskite
solar cells arrive at fundamentally contradictory conclusions regarding
the primary driver of non-radiative recombination and operational
degradation. While one body of literature attributes performance loss at
the buried contact to chemical redox reactions between migrating iodide
species and adjacent charge extraction layers, competing reports argue
that performance decay is governed entirely by field-driven accumulation
of mobile point defects that screen the built-in potential.
Disentangling these mechanisms requires operando, spatially resolved
transient optoelectronic spectroscopy combined with depth-profiled X-ray
photoemission under simultaneous electrical bias and continuous
illumination. This unresolved mechanistic conflict reflects a broader
transition across the discipline. The principal frontier of perovskite
photovoltaics has progressed beyond chasing isolated power conversion
efficiency (PCE) records on sub-\qty{0.1}{\centi\meter\squared}
laboratory apertures. Instead, the central challenges reside in
reproducible third-party verification, maintaining performance across
commercially relevant module areas, and ensuring operating lifetime
under continuous operational stress.

Across the six months from January to June 2026, 3328 indexed works met
the scope gate and 884 were read closely. This body of literature
reveals that closing the gap between small-area laboratory devices and
commercial modules requires resolving several interrelated physical
processes across the device stack. This review follows seven critical
mechanistic threads: trap states, passivation and ionic motion; contact
chemistry at the buried interface; A-site and halide composition
engineering; tandem architecture and current matching; degradation
mechanisms under operation; and large-area coating and the area penalty.
Defect chemistry cannot be isolated from ionic dynamics, as shallow
point defects migrate under operating electric fields, transforming into
deep carrier traps or accumulating at transport layers. Similarly,
contact chemistry at the buried interface dictates initial
crystallization kinetics, internal strain, and energetic alignment,
which dictate long-term resistance to ion migration and delamination.
Meanwhile, A-site and halide composition engineering directly governs
thermodynamic phase stability against segregation under light and
thermal stress, setting the operational boundary for tandem architecture
and current matching in monolithic multijunctions.

Translating these microscopic stabilization strategies to industrially
scalable dimensions exposes the area penalty incurred during large-area
coating, where meniscus drying kinetics, fluid dynamics, and substrate
thermal gradients induce spatial variations in film stoichiometry and
defect density. Under actual operating conditions, these initial
microstructural non-uniformities accelerate localized degradation
mechanisms under operation, precipitating failure through moisture
ingress, heat-induced halide redistribution, and electrode corrosion.
Evaluating these interconnected phenomena across a six-month window
provides a necessary diagnostic perspective: a single month captures
only isolated values and transient records, whereas six months reveal
whether reported performance metrics actually moved or stagnated. A
half-year view cannot resolve multi-thousand-hour degradation
trajectories or guarantee twenty-year service lifetimes, but it offers
sufficient temporal depth to distinguish genuine thermodynamic
stabilization from statistical noise and short-lived laboratory
anomalies.

\subsection{2. The Certified Frontier Across January-June
2026}\label{the-certified-frontier-across-january-june-2026}

A fundamental disagreement governs recent two-terminal (2T) tandem
reports concerning the origin of non-radiative voltage deficits: while
one investigation attributes performance limits to defective
hole-extraction dynamics at the buried self-assembled monolayer (SAM)
contact {[}\href{https://doi.org/10.1126/science.aef5355}{1}{]}, another
concludes that phase segregation and deep trap formation within the bulk
wide-bandgap perovskite dominate carrier recombination
{[}\href{https://doi.org/10.1038/s41467-026-73276-w}{2}{]}. Resolving
this conflict requires operando absolute photoluminescence quantum yield
(PLQY) imaging coupled with transient absorption spectroscopy under
external bias, which can isolate interfacial recombination velocity from
bulk quasi-Fermi level splitting collapse. These opposing mechanistic
attributions yielded certified power conversion efficiency (PCE) values
of 33.6\% for contact optimization
{[}\href{https://doi.org/10.1126/science.aef5355}{1}{]} and 32.95\% for
bulk defect mitigation
{[}\href{https://doi.org/10.1038/s41467-026-73276-w}{2}{]}. Across the
six-month period, independently certified 2T tandem efficiencies
occupied a narrow corridor spanning 31.93\%
{[}\href{https://doi.org/10.1016/j.esci.2026.100563}{3}{]} to 34.85\%
{[}\href{https://doi.org/10.1002/adom.202503341}{4}{]}, illustrating
that divergent physical hypotheses are routinely invoked to explain
marginal variations within the same narrow efficiency envelope.

Over the half year, the best certified PCE moved from 29.2\% in January
to 33.6\% in June, representing a change of 4.4 percentage points.
Rather than demonstrating an advancing technological frontier, this
six-point monthly sequence reflects the stochastic noise of manuscript
indexing schedules. The recorded progression of monthly certified peaks
from 29.2\% in January, to 32.46\% in February
{[}\href{https://doi.org/10.1021/acsenergylett.5c04127}{5}{]}, 31.93\%
in March {[}\href{https://doi.org/10.1016/j.esci.2026.100563}{3}{]},
32.81\% in April {[}\href{https://doi.org/10.1002/adma.73106}{6}{]},
33.5\% in May
{[}\href{https://doi.org/10.1038/s41467-026-72794-x}{7}{]}, and 33.6\%
in June {[}\href{https://doi.org/10.1126/science.aef5355}{1}{]} failed
to establish a newly discovered physical regime. The highest
independently certified efficiency of the entire half year, 34.85\%
{[}\href{https://doi.org/10.1002/adom.202503341}{4}{]}, appeared in
March and exceeded every certified value reported in May and June.
Consequently, the certified ceiling did not move outward, but repeated
within established statistical bounds. Every leading certified result
across this period was achieved on multi-junction devices, specifically
2T tandems, rather than single-junction absorbers. The verified boundary
therefore reflects subcell current matching, optical coupling, and
recombination-junction losses rather than fundamental breakthroughs in
single-junction materials science.

A persistent gap separates uncertified claims from verified physical
metrics: only 114 of 872 closely read papers reported an independently
certified efficiency. In January, 15 papers reported certified values
out of 475 indexed works, with the top self-reported efficiency of
31.85\% standing near the certified 29.2\%. In February, 17 papers
reported certified values out of 499 indexed works, where self-reported
metrics reached 32.73\% alongside an independently certified 32.46\%
{[}\href{https://doi.org/10.1021/acsenergylett.5c04127}{5}{]}. March
included 14 certified papers among 559 indexed works, recording a best
self-reported 32.43\% against a certified 31.93\%
{[}\href{https://doi.org/10.1016/j.esci.2026.100563}{3}{]}. In the
second quarter, uncertified self-reported claims detached completely
from verified performance, escalating to 41.6\% in April across 589
indexed works (21 reporting certified values) and 44.36\% in May across
583 indexed works (29 reporting certified values). By contrast,
independently certified 2T tandem efficiencies converged into a narrow
band, supported by shared interface passivation and additive treatments:
April yielded certified values of 32.81\%
{[}\href{https://doi.org/10.1002/adma.73106}{6}{]} and 32.45\%
{[}\href{https://doi.org/10.1038/s41467-026-72160-x}{8}{]}, while May
produced a dense cluster with 33.5\%
{[}\href{https://doi.org/10.1038/s41467-026-72794-x}{7}{]}, 33.48\%
{[}\href{https://doi.org/10.1126/sciadv.aec4431}{9}{]}, 33.37\%
{[}\href{https://doi.org/10.1002/aenm.71107}{10}{]}, 33.04\%
{[}\href{https://doi.org/10.1038/s41467-026-73656-2}{11}{]}, and 32.1\%
{[}\href{https://doi.org/10.1002/anie.5162616}{12}{]}. June recorded 14
certified reports among 623 indexed works, where self-reported metrics
of 33.59\% fell slightly below certified results at 33.6\%
{[}\href{https://doi.org/10.1126/science.aef5355}{1}{]} and 32.95\%
{[}\href{https://doi.org/10.1038/s41467-026-73276-w}{2}{]}. Independent
certifications clustered tightly across these studies
{[}\href{https://doi.org/10.1038/s41467-026-72794-x}{7},\href{https://doi.org/10.1126/sciadv.aec4431}{9},\href{https://doi.org/10.1002/aenm.71107}{10},\href{https://doi.org/10.1038/s41467-026-73656-2}{11},\href{https://doi.org/10.1002/anie.5162616}{12}{]},
confirming that uncertified claims above 40\% arise from spectral
mismatch errors or omitted aperture masking rather than verified
material mechanisms.

The concentration of high efficiencies in small-area cells exposes the
physical barrier of module scaling, where geometric parasitic losses
depress performance. Across the half year, the share of closely read
papers expanded from January (12.6\%), February (15.0\%), and March
(16.7\%) through April (18.8\%), May (17.7\%), and June (19.2\%),
reflecting increased publication activity around scalable manufacturing.
However, active area expansion introduces severe lateral sheet
resistance in transparent conductive oxide (transparent conducting oxide
(TCO)) electrodes, subcell interconnect dead zones, and non-uniform film
drying, which systematically degrade the fill factor (FF). At
\qty{100}{\centi\meter\squared}, certified module efficiencies were
limited to 26.59\%
{[}\href{https://doi.org/10.1126/sciadv.aec3238}{13}{]} and 24.7\%
{[}\href{https://doi.org/10.1038/s41467-026-69685-6}{14}{]}, where
localized series resistance and laser-scribed interconnect defects
constrain power extraction
{[}\href{https://doi.org/10.1126/sciadv.aec3238}{13},\href{https://doi.org/10.1038/s41467-026-69685-6}{14}{]}.
At larger areas, efficiency diminished further to 27.12\% across
\qty{655.2}{\centi\meter\squared}
{[}\href{https://doi.org/10.1002/adma.73855}{15}{]}, 26.62\% across
\qty{651}{\centi\meter\squared}
{[}\href{https://doi.org/10.1002/aenm.71111}{16}{]}, and 23.3\% across
\qty{694.5}{\centi\meter\squared}
{[}\href{https://doi.org/10.1002/solr.70353}{17}{]}. At the extreme of
reporting, a \qty{900}{\centi\meter\squared} device was presented
without a certified efficiency
{[}\href{https://doi.org/10.1002/adfm.202532205}{18}{]}. These
multi-junction and module data demonstrate that verified small-cell
records above 33\% fail to translate directly into scalable technology,
where macroscopic transport resistance and deposition non-uniformities
supersede nanometer-scale interface engineering.

\subsection{3. Trap States, Passivation and Ionic
Motion}\label{trap-states-passivation-and-ionic-motion}

Recent investigations into trap passivation and ionic transport diverge
sharply regarding the structural configuration required to deactivate
surface defects without impeding charge collection. One faction asserts
that capping the three-dimensional (3D) perovskite absorber with
low-dimensional phases is essential to eliminate residual reagents and
suppress non-radiative recombination. Vapor-mediated dimensional
reconstruction using triethylamine-pentafluorophenol (TEA-PFP) to react
with excess lead iodide (\ce{PbI2}) forms a one-dimensional (1D) TEAPbI3
interfacial phase, yielding a certified power conversion efficiency
(PCE) of 26.15\% (champion 26.71\% on \qty{1}{\centi\meter\squared}) and
a fill factor (FF) of 89.13\%
{[}\href{https://doi.org/10.1038/s41467-026-72790-1}{19}{]}. In situ
generation of a guanidinium derivative similarly templates a
two-dimensional/three-dimensional (2D/3D) heterostructure, giving a
certified PCE of 26.2\% (champion 26.49\%) with 95\% retention after
\qty{1000}{\hour} of maximum power point tracking (MPPT)
{[}\href{https://doi.org/10.1002/anie.8890696}{20}{]}. Conversely,
competing reports argue that low-dimensional capping phases introduce
anisotropic transport barriers, carrier accumulation, and phase
instability under operational bias. They demonstrate instead that non-2D
organic co-crystal layers passivate defects via reinforced
hydrogen-bonding networks and enhanced electron mobility, achieving a
champion PCE of 26.35\% with a T80 lifetime of \qty{1392}{\hour} and
88\% efficiency retention after \qty{2040}{\hour} at approximately 50\%
relative humidity
{[}\href{https://doi.org/10.1002/adma.202519138}{21}{]}. This rejection
of low-dimensional capping is corroborated by stereoelectronic ligand
engineering: pyridine or pyrimidine rings forming dual lead-nitrogen
(Pb-N) coordination and Pb-I-pi interactions drive planar ligand
alignment, producing atomic-scale defect mitigation with sub-nanometre
charge transfer and a certified PCE of 27.41\% (stabilised at 26.85\%)
without dimensional reconstruction
{[}\href{https://doi.org/10.1038/s41586-026-10626-0}{22}{]}. Settling
this dispute requires operando bias-dependent, depth-resolved
grazing-incidence wide-angle X-ray scattering paired with local
transient absorption microscopy under continuous thermal stress to
establish whether 1D and 2D phases degrade into resistive barriers under
operational electric fields or preserve a stable passivation contact.

Chemical passivation strategies have targeted undercoordinated metal
cations, halide vacancies, and antisite defects through coordinated
Lewis base and acid interactions. Carboxyl-group chelation of
undercoordinated lead (\ce{Pb^2+}) by dipicolinic acid (DPA) passivates
surface and bulk traps in \qty{1.68}{\electronvolt} wide-bandgap (WBG)
cells, reaching 22.91\% PCE and a champion 31.66\% in two-terminal (2T)
perovskite/silicon tandems
{[}\href{https://doi.org/10.1002/aenm.70902}{23}{]}. Lewis acid-base
coordination using 3-hydroxyflavone coordinates \ce{Pb^2+} and increases
the ion-migration barrier, though its champion PCE of 26.6\% was
measured on an uncertified, small area of
\qty{0.09}{\centi\meter\squared}
{[}\href{https://doi.org/10.1002/smll.73673}{24}{]}. Adding
meta-fluoro-phenylethylammonium iodide (m-F-PEAI) to ethyl acetate
antisolvent employs the fluorine atom as a Lewis base toward \ce{Pb^2+}
while establishing directed hydrogen bonds with formamidinium iodide
(FAI) to retard crystallization, achieving 26.41\% champion PCE
{[}\href{https://doi.org/10.1021/acsaem.6c00562}{25}{]}. Redox-active
dabrafenib (DRFN) coordinates multiple \ce{Pb^2+} centers via its
sulfonyl group and heterocyclic nitrogen, achieving peak PCEs of 25.28\%
in n-i-p and 26.62\% in p-i-n devices
{[}\href{https://doi.org/10.1002/adfm.76623}{26}{]}. Molecular dipole
passivators 2-amino-4,6-dichlorotriazine and 1-adamantanecarbonyl
chloride hydrochloride (ACHCl) combine Lewis-base coordination with
chloride filling of iodine vacancies (VI), correcting PbI6 octahedral
distortion and reducing recombination to reach champion PCEs of 26.58\%
and 26.12\%, respectively
{[}\href{https://doi.org/10.1021/acsmaterialslett.5c01678}{27},\href{https://doi.org/10.1007/s40820-026-02150-7}{28}{]}.
Multi-defect healing additives pairing pentafluorophenol (PFP) with
1,4-dicyanobenzene (DCB) target VI, lead vacancies (VPb), and
lead-on-iodine (PbI) antisites, retaining 95\% efficiency after
\qty{1300}{\hour} at 65 °C and 94\% after \qty{1000}{\hour} of 1-sun
MPPT (champion PCE 26.47\%)
{[}\href{https://doi.org/10.1002/smll.73436}{29}{]}. Complementary
bimolecular passivation with 1,4-phenylenediammonium diiodide (PDADI)
and 11-mercaptoundecanoic acid (MUA) reduces trap density from 3.00 x
10\^{}15 to 1.65 x 10\textsuperscript{\qty{15}{\centi\meter}}-3
(champion PCE 26.15\%)
{[}\href{https://doi.org/10.1002/adfm.202532110}{30}{]}, while
quantitative inkjet printing of PDADI alone yields 26.1\% PCE
{[}\href{https://doi.org/10.1002/smll.73886}{31}{]}. In the bulk,
precursor dicyandiamide improves vacancy formation energies and
crystallinity, producing a monolithic perovskite/silicon tandem champion
PCE of 32.03\% with an open-circuit voltage (Voc) of \qty{1.95}{\volt}
{[}\href{https://doi.org/10.1021/acsami.6c01711}{32}{]}. Chiral
configurations show that stereochemical orientation controls defect
thermodynamics: (R)-(+)-alpha-phenylethylamine (R-Pye) exhibits higher
binding energy (9.45 versus \qty{7.33}{\electronvolt}) and increased VI
formation energy (4.56 versus \qty{2.46}{\electronvolt}) than
(S)-(-)-alpha-phenylethylamine (S-Pye), reaching 26.13\% champion PCE
{[}\href{https://doi.org/10.1002/aenm.71222}{33}{]}. To prevent ammonium
deprotonation, sterically protected cations like
N,N-dimethylcyclohexylammonium iodide (DMCHAI) inhibit side reactions
with formamidinium (FA+), maintaining over 86\% efficiency after
\qty{550}{\hour} at 85 °C (champion PCE 26.06\%)
{[}\href{https://doi.org/10.1021/acsami.6c04956}{34}{]}, matching steric
and electronic tuning in trimetazidine (TMZ) that achieves 26.07\% PCE
with over 90\% retention after \qty{1400}{\hour}
{[}\href{https://doi.org/10.1002/adfm.75743}{35}{]}. Latent passivator
3-trifluoromethylphenylboronic acid pinacol ester (3-CF3-PhBpin)
hydrolyzes in situ into 3-CF3-PhB(OH)2 to passivate crystallization
defects, achieving 26.25\% in rigid and 24.80\% in flexible n-i-p cells
{[}\href{https://doi.org/10.1002/adfm.75370}{36}{]}. Machine-learning
frameworks integrating domain-knowledge large language models with
Bayesian optimization have independently proposed untested
four-component recipes (3MTPAI, PDAI2, EDAI2, and PipDI) reaching 26.0\%
champion PCE {[}\href{https://doi.org/10.48550/arxiv.2602.04914}{37}{]}.

Over the six-month period, investigation of ion migration and halide
segregation revealed an incompatibility between static immobilization
and dynamic strain dissipation. Static strategies pin mobile ions
through spatial confinement or covalent polymerization. Molecular
confinement domains suppress halide segregation by constraining FA+
motion, enabling single-junction efficiencies of 24.21\%
(\qty{1.68}{\electronvolt}) and 21.20\% (\qty{1.78}{\electronvolt}) and
a 33.59\% champion 2T tandem PCE
{[}\href{https://doi.org/10.1002/adma.73713}{38}{]}. Incorporating
phenylethylguanidinium (PEGA+) salts suppresses cation drift in
\qty{1.55}{\electronvolt} perovskites, exceeding 26.2\% PCE with over
96\% retention after \qty{1100}{\hour} and 89\% after \qty{800}{\hour}
at 85 °C {[}\href{https://doi.org/10.1021/acsenergylett.6c00737}{39}{]}.
In situ polymerization of vinylphosphonic acid outclasses small-molecule
organoammoniums, retaining over 90\% PCE after \qty{1600}{\hour} of
1-sun MPPT (certified 26.24\%, champion 26.54\%)
{[}\href{https://doi.org/10.1002/adma.202522825}{40}{]}. In situ
polymerization of sodium thioctate (TANA) combined with sodium
interstitial doping regulates crystallization and suppresses ion
migration, retaining 92.3\% efficiency after \qty{1000}{\hour} under
International Summit on Organic Photovoltaic Stability (ISOS) ISOS-L-2I
conditions (certified 26.75\%, champion 26.9\% on a
\qty{10.24}{\centi\meter\squared} module)
{[}\href{https://doi.org/10.1002/anie.202525815}{41}{]}.
Self-polymerizable fluorinated additives (PFPA) achieve a champion PCE
of 26.01\% on a \qty{14}{\centi\meter\squared} area, retaining 90.0\%
efficiency after \qty{840}{\hour} at 85 °C in nitrogen and 95.1\% after
\qty{1000}{\hour} of MPPT
{[}\href{https://doi.org/10.1021/acsenergylett.5c03901}{42}{]}.
Room-temperature condensation between amino and acyl chloride groups
establishes a dual-anchor polymeric interlayer, preserving 85\%
efficiency after \qty{2300}{\hour} of MPPT and \qty{1680}{\hour} at 65
°C (champion PCE 26.12\%)
{[}\href{https://doi.org/10.1002/anie.7579788}{43}{]}. At the transport
interface, bisphosphonate interlayers (2OBCz) between the perovskite and
hole transport layer (hole-transport layer (HTL)) inhibit ion migration,
retaining 90.7\% efficiency after \qty{1000}{\hour} under ISOS-D-2
conditions compared to monophosphonates (LBCz) (champion PCE 26.02\%)
{[}\href{https://doi.org/10.1002/adfm.75412}{44}{]}. In narrow-bandgap
tin-lead (Sn-Pb) cells, formamidine bisulfate (FBH) dipole management
reverses top-surface grain-boundary states to induce downward band
bending, preserving over 80\% efficiency after \qty{400}{\hour} at 85 °C
and 90\% after \qty{1872}{\hour} in a glovebox (champion PCE 23.25\%,
tandem 29.67\%) {[}\href{https://doi.org/10.1002/adma.72755}{45}{]}.

In direct conflict with the static immobilization premise, evidence
indicates that rigid polymer networks suffer fatigue and microcracking
under thermal cycling, opening pathways for ion accumulation. Dynamic
networks resolve this failure mode: mechanochemically reinforced
dual-dynamic covalent seeding (DDS) via orthogonal Diels-Alder and
oxime-carbamate reactions constructs an interpenetrating covalent mesh
at grain boundaries that heals under strain while passivating defects
and blocking ion and moisture ingress (certified PCE 26.69\%, champion
26.95\%, ISOS-L-2I) {[}\href{https://doi.org/10.1002/adma.72823}{46}{]}.
Integrating photoswitchable isomers at grain boundaries suppresses
lattice bond rupture and defect accumulation during cyclic illumination
and thermal fluctuation, retaining over 95\% of initial efficiency after
2,\qty{000}{\hour} of UV-containing light cycling at 65 °C and 500
thermal cycles between -40 °C and 85 °C (certified PCE 26.9\%, champion
27.2\%) {[}\href{https://doi.org/10.1038/s41560-026-01993-z}{47}{]}.
Dynamic methyl viologen (MV) interlayers similarly maintain 90\%
efficiency after 1,\qty{000}{\hour} of AM 1.5G MPPT (champion PCE
26.44\%, Voc \qty{1.167}{\volt}, FF 85\%)
{[}\href{https://doi.org/10.1002/adfm.202530317}{48}{]}. Systematic
temperature-dependent transient ion drift measurements across sustained
thermal cycles are necessary to determine whether dynamic covalent
networks surpass rigid polymers in preventing long-term mobile ion
accumulation.

Surface passivation and ionic stabilization depend heavily on energetic
alignment across selective contacts. In inverted p-i-n devices,
modulating intermolecular N-H\ldots N hydrogen bonding between two
passivators tunes energy-level offsets, yielding a 24.04\% wide-bandgap
single-junction PCE and a certified monolithic 2T tandem PCE of 33.48\%
(champion 33.63\% on \qty{1}{\centi\meter\squared})
{[}\href{https://doi.org/10.1126/sciadv.aec4431}{9}{]}. Self-assembled
monolayers (SAMs) on nickel oxide (\ce{NiO_x}) suffer from phase
aggregation and incomplete coverage. Co-adsorbing N-indoleacetic acid
with {[}4-(3,6-dimethyl-9H-carbazol-9-yl)butyl{]}phosphonic acid
(Ph-4PACz) forms phase-homogeneous binary monolayers via pi-pi stacking
and hydrogen bonding, achieving 26.95\% PCE (certified 26.57\%) on
\qty{0.0717}{\centi\meter\squared} and 25.61\% on
\qty{1}{\centi\meter\squared}
{[}\href{https://doi.org/10.1038/s41467-026-71301-6}{49}{]}. Blending
4PABCz with Me-4PACz forms a Co-SAM that enhances coverage on
\ce{NiO_x}, passivating bottom defects and suppressing non-radiative
recombination (champion PCE 26.33\%)
{[}\href{https://doi.org/10.1002/eem2.70443}{50}{]}. Positional
isomerism provides structural control: selective chlorine placement on
the peripheral phenyl ring of Ph-4PACz (ortho-, meta-, and
para-chlorinated isomers) enables lattice-matched dual-site passivation
(certified PCE 26.13\%, champion 26.58\%)
{[}\href{https://doi.org/10.1002/adma.202523249}{51}{]}.
Indolocarbazole-based SAMs (M3PAICz-1) demonstrate versatility across
bandgaps, delivering 26.12\% PCE in \qty{1.55}{\electronvolt} cells and
22.19\% in \qty{1.68}{\electronvolt} wide-bandgap devices
{[}\href{https://doi.org/10.1002/smll.202512942}{52}{]}, while tripodal
phosphonic acids anchor multidentately with face-on pi-stacking and
segregate to grain boundaries to passivate undercoordinated \ce{Pb^2+}
(certified PCE 26.35\%)
{[}\href{https://doi.org/10.1002/aenm.70806}{53}{]}. The underlying
metal-oxide contact can also be chemically stabilized: bromoacetamide
(BAA) undergoes in situ dissociative adsorption on sputtered \ce{NiO_x},
where bromide ions coordinate surface nickel to stabilize \ce{Ni^3+} and
suppress disorder (certified PCE 26.07\%, champion 26.31\%)
{[}\href{https://doi.org/10.1002/adma.73382}{54}{]}.

At the electron-selective contact, isomer-free multi-adduct fullerene
derivatives \ce{C60}(NHR)4O (TATPC) blended into
{[}\href{https://doi.org/10.1002/adma.73106}{6}{]}-phenyl-C61-butyric
acid methyl ester (PCBM) raise the lowest unoccupied molecular orbital
(LUMO) level, strengthening interfacial interactions to reach 26.66\%
champion PCE on a \qty{1.04}{\centi\meter\squared} aperture
{[}\href{https://doi.org/10.1002/anie.4500167}{55}{]}. Fused-ring
electron acceptor (FREA) additives similarly passivate defects and
accelerate electron extraction, yielding a certified PCE of 26.35\%
(champion 26.66\%) with under 8\% degradation after \qty{2150}{\hour}
under continuous illumination
{[}\href{https://doi.org/10.1021/jacs.5c16976}{56}{]}. In tandem
subcells, differential coordination using 1,10-phenanthroline
coordinates \ce{Sn^2+} more strongly than \ce{Pb^2+}, synchronizing
crystallization to deliver a Sn-Pb single-junction Voc of
\qty{0.887}{\volt}, an FF of 81.27\%, and a PCE of 23.63\% (tandem
champion 29.08\%) {[}\href{https://doi.org/10.1002/eem2.70453}{57}{]}.
Dual-molecule precursor additives (AGS and AHC) similarly coordinate
Sn-Pb crystallization, achieving 23.44\% single-junction efficiency and
a certified tandem PCE of 28.57\% (champion 29.36\%)
{[}\href{https://doi.org/10.1002/adma.73038}{58}{]}, while halide-cation
interactions tuned via chlorine and fluorine salts yield wide-bandgap
subcells of 20.5\% and 20.3\% PCE, retaining 87\% tandem efficiency
after \qty{700}{\hour} of MPPT (tandem champion 29.3\%)
{[}\href{https://doi.org/10.1002/adma.202522668}{59}{]}. Interface
passivation extends to the silicon bottom cell; dark annealing of n-type
tunnel oxide passivated contact (n-TOPCon) silicon cells at 300 °C
reduces defect-assisted recombination, raising subcell efficiency from
18.92\% to 20.23\% in four-terminal tandems with a 26.75\% champion PCE
{[}\href{https://doi.org/10.1002/ente.70501}{60}{]}. In comparison,
idealized one-dimensional numerical simulations (SCAPS-1D) for AgCdF3
absorbers with transition-metal dichalcogenide transport layers
demonstrate that reducing bulk defect density to
\qty{1e13}{\per\centi\meter\cubed} at \qty{0.4}{\micro\meter} thickness
defines an efficiency benchmark of 31.92\%
{[}\href{https://doi.org/10.1002/ente.70553}{61}{]}.

\begin{figure}[htbp]
\centering
\includegraphics[width=0.95\textwidth]{fig/F3_effort_map.pdf}
\caption{Where the month's mechanistic effort was directed. Depth-tier papers by mechanism axis, stacked by device family.}
\end{figure}

\subsection{4. Contact Chemistry at the Buried
Interface}\label{contact-chemistry-at-the-buried-interface}

A fundamental contradiction defines recent investigations into carrier
transfer across self-assembled monolayer (SAM) contacts at the buried
interface of perovskite solar cells (PSCs). On one side, prevailing
models treat hole extraction as a static electrostatic problem governed
by ground-state vacuum work functions and interfacial dipole alignment
{[}\href{https://doi.org/10.1002/adfm.76291}{62},\href{https://doi.org/10.1002/adma.202523338}{63},\href{https://doi.org/10.1002/anie.8473744}{64}{]}.
In this framework, introducing polar co-adsorbents such as cyanoacetic
acid (CA) to raise the work function of
{[}4-(3,6-dimethyl-9H-carbazol-9-yl)butyl{]}phosphonic acid (Me-4PACz)
yields an uncertified champion power conversion efficiency (PCE) of
26.52\% {[}\href{https://doi.org/10.1002/adfm.76291}{62}{]}, while
assembling asymmetric Br-4PhCz monolayers with large dipole moments
produces an open-circuit voltage (\(V_\mathrm{oc}\)) of
\qty{1.36}{\volt} in \qty{1.77}{\electronvolt} wide-bandgap cells,
delivering a certified tandem PCE of 28.98\% (champion 29.35\%) on a
small \qty{0.05}{\centi\meter\squared} area
{[}\href{https://doi.org/10.1002/adma.202523338}{63}{]}. Similarly,
terminal group tuning alone is reported to suppress non-radiative
recombination to reach 95\% and 91\% of the Shockley-Queisser limit in
\qty{1.53}{\electronvolt} (uncertified champion PCE of 26.88\%) and
\qty{1.72}{\electronvolt} devices
{[}\href{https://doi.org/10.1002/anie.8473744}{64}{]}. Conversely,
direct operational interrogation reveals that charge transfer across
SAM-modified interfaces is dynamic and governed by a self-limiting
negative feedback loop
{[}\href{https://doi.org/10.1002/adfm.75295}{65}{]}. In this opposing
mechanism, photo-injected holes accumulating at the buried contact
establish a transient opposing electric field that induces rapid field
saturation over short characteristic times, triggering a sharp decay in
extraction efficacy regardless of favorable initial band offsets
{[}\href{https://doi.org/10.1002/adfm.75295}{65}{]}. These two positions
cannot be reconciled: if transient space-charge accumulation dominates
working conditions, static energy level optimization cannot prevent
extraction decline under continuous illumination. Resolving this
conflict requires operando, time-resolved Stark-shift spectroscopy or
bias-dependent transient electro-absorption measurements under
continuous illumination to determine whether static interfacial dipoles
persist or undergo severe electrostatic screening during steady-state
carrier flux.

Beyond dipolar tuning, modifying electronic coupling across buried hole
transport layers (HTLs) has targeted both chemical linker architecture
and non-radiative decay channels. Inserting thiophene rings into
sulfur-modulated, fully conjugated monolayers (SMe-TPA-ThPA and
SMe-TPA-PA) raises inverted device efficiencies from 24.31\% in
oxygen-linked controls to 26.52\% and 25.40\% by strengthening
intermolecular pi-electron interactions and energetic alignment
{[}\href{https://doi.org/10.1002/anie.5319815}{66}{]}. Interfacial
polarization can also be tuned through meta-substituted bifunctional
molecules such as m-APCA, which suppresses SAM aggregation to achieve a
certified 26.71\% PCE (champion 26.77\%)
{[}\href{https://doi.org/10.1002/adma.73306}{67}{]}, or by introducing a
4-aminobutylphosphonic acid (4-ABPA) molecular bridge that limits
interfacial voltage loss to \qty{31}{\milli\volt}, yielding uncertified
efficiencies of 25.56\% in n-i-p and 26.45\% in p-i-n configurations
{[}\href{https://doi.org/10.1002/adma.202519267}{68}{]}. Rather than
relying solely on electrostatic barrier adjustment, resonant energy
transfer channels have been established: incorporating the UV-312
molecule introduces a Förster resonance energy transfer (FRET) pathway
that facilitates charge separation while suppressing ultraviolet-induced
lead-iodide bond dissociation, achieving an uncertified 27.05\% PCE with
a \(V_\mathrm{oc}\) of \qty{1.186}{\volt} on a
\qty{0.09}{\centi\meter\squared} aperture
{[}\href{https://doi.org/10.1002/anie.3596437}{69}{]}. However, the
efficacy of these small-molecule modifiers depends entirely on substrate
anchoring density and the absence of microscopic pinholes.

The second physical problem centers on anchoring density, molecular
orientation, and phase segregation on rough or chemically heterogeneous
substrates. Small-molecule phosphonic acids bind erratically to
transparent conducting oxides because of non-uniform hydroxyl
distributions. Over the past six months, synthetic strategies shifted
from modifying underlying oxides toward catalytic assembly and
macromolecular networks. Initial interventions altered substrate surface
chemistry directly: depositing three atomic layers of aluminum oxide via
atomic layer deposition (ALD) generates dense hydroxyl sites that align
self-assembled monolayers (SAMs) vertically, yielding an uncertified
26.63\% PCE on \qty{0.09}{\centi\meter\squared} and 22.97\% on a
\qty{14.40}{\centi\meter\squared} mini-module
{[}\href{https://doi.org/10.1021/acsami.6c03372}{70}{]}. In two-terminal
perovskite/silicon tandems, incorporating carbon into polysilicon tunnel
recombination layers elevates surface hydroxyl density to anchor SAMs
securely, producing a certified 33.5\% tandem PCE (champion 33.84\%) on
\qty{1}{\centi\meter\squared} with a T80 lifetime of \qty{800}{\hour}
{[}\href{https://doi.org/10.1038/s41467-026-72794-x}{7}{]}. Similarly,
an annealed Ti3C2Tx MXene interlayer forms tridentate Ti-O-P bonds with
phosphonic acid heads on indium tin oxide (ITO), attaining a certified
26.44\% PCE (champion 26.78\%) on a \qty{0.08}{\centi\meter\squared}
area {[}\href{https://doi.org/10.1002/aenm.70966}{71}{]}. Bypassing
substrate pre-treatments, introducing 4-dimethylaminopyridine (DMAP) as
a solution catalyst accelerates phosphonic acid condensation via a
transient phosphorylpyridinium intermediate, densifying the SAM to
achieve a certified 26.44\% PCE (champion 26.93\%) with 95\% retention
after \qty{700}{\hour}
{[}\href{https://doi.org/10.1021/acsenergylett.6c01257}{72}{]}. Where
pinholes persist on nickel oxide (\ce{NiO_x}), area-selective ALD of
AlOx passivates exposed oxide regions without coating SAM-covered
domains, enabling a \qty{64.68}{\centi\meter\squared} module PCE of
20.88\% (champion 26.41\%)
{[}\href{https://doi.org/10.1002/anie.202516537}{73}{]}, while treating
\ce{NiO_x} with cesium oxalate (CsOA) converts \ce{Ni^3+} into
{[}Ni(C2O4)3{]}3- and eliminates \ce{Ni^4+} sites, supporting an
\qty{11.0}{\centi\meter\squared} device PCE of 26.48\% with 85.7\%
retention after \qty{1560}{\hour} at 65 °C
{[}\href{https://doi.org/10.1002/anie.4012708}{74}{]}.

Molecular shape has likewise been manipulated to regulate self-assembly
thermodynamics. Incorporating dimethyl substituents onto spirofluorene
cores imposes steric hindrance to force ordered molecular packing,
yielding an uncertified champion PCE of 26.62\% that maintains
approximately 90\% of its initial value after \qty{1000}{\hour} of
maximum power point tracking (MPPT) in nitrogen
{[}\href{https://doi.org/10.1002/aenm.71005}{75}{]}. Bisphosphonate
anchors in MeS-BPADCP enforce vertical orientation through bidentate
binding, reaching an uncertified 26.73\% PCE with 94.2\% retention after
\qty{1000}{\hour} under one-sun illumination
{[}\href{https://doi.org/10.1002/aenm.70814}{76}{]}. A butterfly-shaped
dibenzazepine core (PATDBAz) enables pi-pi stacking while preventing
aggregation, sustaining 97.5\% of initial efficiency after
\qty{1000}{\hour} under International Summit on Organic Photovoltaic
Stability (ISOS) ISOS-L-1 testing (champion 26.47\%)
{[}\href{https://doi.org/10.1002/smll.74278}{77}{]}. Intermolecular
ordering can also be modulated by doping mixed SAMs to alleviate tensile
interfacial strain, achieving a tandem champion PCE of 30.18\% and an
inverted \qty{1.68}{\electronvolt} cell PCE of 23.48\% (\textgreater90\%
retention after \qty{1500}{\hour})
{[}\href{https://doi.org/10.1002/smll.74319}{78}{]}, or by
post-assembling MeO-ADCP onto pre-formed RS-2 monolayers to fill bare
vacancies, achieving an uncertified 26.8\% PCE with over 99\% retention
after \qty{2000}{\hour} at 45 °C
{[}\href{https://doi.org/10.1021/acs.jpclett.6c01042}{79}{]}. Despite
these ordering schemes, monomeric SAMs exhibit weak shear resistance and
detach under thermomechanical stress. Consequently, research shifted
toward polymeric networks. Replacing small molecules with an asymmetric
polycarbazole phosphonic acid (Poly-3Ph-4PACz) homogenizes flexible
hole-selective contacts, yielding a certified tandem PCE of 32.95\%
(champion 26.13\%, stabilised 27.12\%) and a
\qty{57.6}{\centi\meter\squared} flexible module PCE of 22.22\%
{[}\href{https://doi.org/10.1038/s41467-026-73276-w}{2}{]}. Multidentate
phosphonic acid polymer networks (P-SAM) increase interfacial fracture
toughness nearly eightfold to 9.11 MPa (certified 26.61\% on
\qty{0.06734}{\centi\meter\squared})
{[}\href{https://doi.org/10.1002/aenm.71159}{80}{]}, while the
phosphonic-acid-functionalized polymer PTPP combines continuous film
formation with anchoring stability, sustaining 94.4\% of initial
performance after 800 thermal cycles (certified 26.4\%, champion 27.1\%)
{[}\href{https://doi.org/10.1002/adfm.75697}{81}{]}. In tin-lead
absorbers, blending hydrophobic trifluoromethyl molecules (MTFP) into
poly(3,4-ethylenedioxythiophene):poly(styrenesulfonate) (PEDOT:PSS)
suppresses moisture uptake to reach 23.36\% PCE (\(V_\mathrm{oc}\) of
\qty{0.88}{\volt}) and a tandem PCE of 28.67\%
{[}\href{https://doi.org/10.1002/aenm.70910}{82}{]}.

The third physical problem involves chemical coordination across the
buried interface, which dictates heterogeneous nucleation, trap
passivation, and ion migration. Buried contact chemistry directly
regulates crystallization kinetics: imidazolium-functionalized SAMs
guide crystal growth through pi-cooperative coordination with lead
polyhalide complexes, retaining 93.8\% of initial PCE after
\qty{936}{\hour} of operation at the maximum power point (champion
26.83\%) {[}\href{https://doi.org/10.1002/adma.73828}{83}{]}. On
flexible substrates where poor precursor wetting impedes grain
formation, modifying the SAM with NCGCl salt utilizes nitrile and
guanidinium groups to promote heterogeneous nucleation and passivate
defects, achieving a certified 26.64\% PCE (champion 26.89\%)
{[}\href{https://doi.org/10.1002/adma.73358}{84}{]}. Precursor
pre-organization before coating also restricts phase competition: the
zwitterionic additive CPTAC associates with formamidinium cations,
contracting competing crystallization pathways to accelerate conversion
into the photoactive alpha phase (uncertified 26.54\% PCE on
\qty{0.1}{\centi\meter\squared})
{[}\href{https://doi.org/10.1021/acsenergylett.6c01291}{85}{]}.

Passivating deep defects at the buried interface has converged on
chelating agents, heterojunction formation, and diffusion barriers.
Inserting an aromatic amino acid derivative (NBP) between Me-4PACz and
wide-bandgap absorbers coordinates unbonded \ce{Pb^2+} cations and
passivates iodide vacancies, enabling wide-bandgap single junctions to
reach 20.81\% PCE with a \(V_\mathrm{oc}\) of \qty{1.371}{\volt} and
tandems to reach 29.05\%
{[}\href{https://doi.org/10.1002/smll.73595}{86}{]}. Bifunctional
cyanoacetic acid incorporated into both the SAM and perovskite precursor
passivates interface boundaries, suppresses ion migration, and nearly
doubles the reverse-bias breakdown voltage, achieving a tandem champion
of 32.43\% (stabilised 24.43\%)
{[}\href{https://doi.org/10.1002/adma.72738}{87}{]}. In situ growth of
gradient Ruddlesden-Popper two-dimensional/three-dimensional (2D/3D)
heterostructures suppresses interface defects, reducing tandem current
mismatch from 5.5\% to 0.1\% (certified 32.46\%, champion 32.73\%)
{[}\href{https://doi.org/10.1021/acsenergylett.5c04127}{5}{]}, while
conjugated TzTzI ligands construct 2D/3D networks that yield 24.5\% on
\qty{1}{\centi\meter\squared}, 20.61\% for \qty{1.72}{\electronvolt}
bandgaps, and 22.14\% in mini-modules (champion 26.61\%)
{[}\href{https://doi.org/10.1002/adma.202520251}{88}{]}. Multisite
defect healing with thiazolidine derivatives (TZCl-CN) suppresses
non-radiative recombination to reach an uncertified 26.39\% PCE
{[}\href{https://doi.org/10.1021/acsenergylett.5c03328}{89}{]}.
Bypassing HTLs entirely, bimolecular co-modification with THDPH and
PEABr manages the perovskite/electrode boundary, retaining 90\% of
initial efficiency after \qty{1000}{\hour} (champion 26.55\%)
{[}\href{https://doi.org/10.1002/adfm.76198}{90}{]}. To prevent lead
leakage, polymerizing the crown ether DVB18C6 in situ yields a 95.14\%
\ce{Pb^2+} sequestration efficiency (champion 26.78\%)
{[}\href{https://doi.org/10.1021/acsenergylett.6c00703}{91}{]}.
Ultrathin inorganic interlayers provide physical barriers to ion
migration: ALD HfOx layers anchor 3F-PEAI via Hf-F interactions to block
halide and silver diffusion (certified 26.6\%, champion 27.1\%)
{[}\href{https://doi.org/10.1126/science.aea3339}{92}{]}, paralleling
SAM architectures that lose only 2\% efficiency after \qty{2100}{\hour}
of MPPT at 65 °C under ISOS-L-2 conditions (certified 26.90\%, champion
27.10\%) {[}\href{https://doi.org/10.1038/s41467-026-73426-0}{93}{]}.

Buried electron transport layers (ETLs) present analogous chemical and
electrostatic challenges. Dual passivation combining DMDPTF with
cobalt-acetate-modified tin oxide (\ce{SnO2}) produces an uncertified
26.41\% PCE on \qty{0.096}{\centi\meter\squared} and 22.18\% on a
\qty{10.04}{\centi\meter\squared} module
{[}\href{https://doi.org/10.1002/smll.73891}{94}{]}. Reconstructing
\ce{SnO2} nanocrystals with in situ generated {[}Al(OH)4{]}- ions
rebuilds the interfacial electrical double layer, delivering a certified
26.64\% PCE (champion 26.7\%) and 21.91\% on a
\qty{65}{\centi\meter\squared} module
{[}\href{https://doi.org/10.1002/adma.73286}{95}{]}. For indoor
applications, Mg-doping of \ce{SnO2} widens the optical bandgap and
reduces recombination to yield 35.54\% under light-emitting diode (LED)
illumination and 20.28\% under one sun
{[}\href{https://doi.org/10.1002/smll.202512808}{96}{]}, while dual
insulator contacts block shunting pathways to deliver an indoor PCE of
44.36\% at 1000 lx and 43.08\% at 600 lx (\(V_\mathrm{oc}\) of
\qty{1.064}{\volt}, fill factor of 84.52\%)
{[}\href{https://doi.org/10.1007/s40820-026-02225-5}{97}{]}.

\subsection{5. A-Site and Halide Composition
Engineering}\label{a-site-and-halide-composition-engineering}

Investigations into formamidinium lead triiodide (\ce{FAPbI3}) phase
stabilization reach conflicting conclusions on how low-dimensional
templates suppress the alpha-to-delta transition. Contact-triggered
cationic interaction (CCI) asserts that structural refinement, cation
homogenization, and defect suppression proceed via non-destructive
physical contact between separately crystallized two-dimensional (2D)
and three-dimensional (3D) perovskites without chemical mixing,
interdiffusion, or permanent bonding (certified 25.61\% power conversion
efficiency (PCE), champion 26.25\%)
{[}\href{https://doi.org/10.1038/s41560-026-02027-4}{98}{]}. Conversely,
formamidinium tetrafluoroborate (FABF4) recrystallization demonstrates
that constructing stable mixed-dimensional networks requires active
formamidinium (FA+)/methylammonium (MA+) cation exchange and Ostwald
ripening rather than passive contact, achieving an uncertified 21.36\%
champion PCE that retains 88\% of initial efficiency after 10,000
bending cycles in nitrogen and 82\% at 60\% ± 5\% relative humidity (RH)
{[}\href{https://doi.org/10.1002/adfm.75673}{99}{]}. This conflicts with
crystallographic interface models where the alpha-to-delta transition
nucleates epitaxially at a tensilely strained, coherent interface
between alpha-phase (011) and delta-phase (110) planes, demanding 2D
perovskitoid interfaces to disrupt delta-phase orientation (uncertified
champion 25.61\% PCE)
{[}\href{https://doi.org/10.1038/s41467-026-73739-0}{100}{]}. Resolving
whether stabilization requires passive contact, active cation exchange,
or orientation disruption demands operando cross-sectional depth
profiling (such as cryogenic secondary ion mass spectrometry) paired
with transmission electron microscopy under thermal bias to establish
whether bulky A-site cations stay segregated or interdiffuse across the
interface.

Precursor stabilization strategies over the past six months divide over
whether failure originates from proton transfer or physical volatility.
Competitive proton-transfer suppression using methyl trifluoroacetate
prevents FA+ deprotonation to extend precursor shelf life beyond 22
days, producing uncertified champion devices reaching 26.35\% PCE that
retain 94.75\% efficiency after \qty{1464}{\hour} of humidity exposure
{[}\href{https://doi.org/10.1038/s41467-026-73620-0}{101}{]}. In single
crystals, citric acid similarly mitigates cation deprotonation,
scavenges oxidized iodine (I2/I3-), and passivates lead (\ce{Pb^2+})
dangling bonds to yield 23.39\% PCE
{[}\href{https://doi.org/10.1002/aenm.70939}{102}{]}, while multidentate
ZL1 anchors suppress volatile FA loss to lift PCE from 24.20\% to
26.13\% {[}\href{https://doi.org/10.1007/s40820-026-02098-8}{103}{]}. In
contrast, studies on volatile additives demonstrate that cation
volatility dictates structural evolution: in thiocyanate systems,
ammonium (\ce{NH4}+) evaporates completely during annealing to drive
bulk crystallization, whereas non-volatile phenylethylammonium (PEA+)
segregates to grain boundaries for interfacial passivation (23.3\% PCE)
{[}\href{https://doi.org/10.1021/acsaem.6c01329}{104}{]}. In cesium lead
triiodide (\ce{CsPbI3}), formamidinium acetate (FAAC) intermediate
engineering lowers Cs4PbI6 formation energy, where transient FA+
intercalation into DMAPbI3 accelerates dimethylammonium (DMA+)/cesium
(\ce{Cs+}) exchange and dimethylammonium iodide (DMAI) volatilization
(21.84\% PCE)
{[}\href{https://doi.org/10.1021/acs.jpclett.6c01395}{105}{]}.
Conversely, ethyl thioacetate post-treatment bypasses cation exchange to
passivate uncoordinated \ce{Pb^2+} through Lewis base coordination
(17.07\% PCE; T80 of \qty{1360}{\hour})
{[}\href{https://doi.org/10.1021/acsaem.5c04098}{106}{]}. Rubidium
chloride (RbCl) guides phase transformations in printable mesoscopic
methylammonium lead triiodide (\ce{MAPbI3}), though performance is
constrained to 18.48\% PCE with only \qty{300}{\second} of steady-state
tracking {[}\href{https://doi.org/10.1002/solr.70363}{107}{]}. Machine
learning via LEAP couples large language models with active learning to
prioritize 6-CDQ and 2-CNA additives, lifting average PCE from a 19.25\%
control to 20.13\% and 20.87\% (21.32\% champion)
{[}\href{https://doi.org/10.48550/arxiv.2605.20242}{108}{]}, while
vacuum-assisted urea coordination modulates triple-cation
crystallization to reach 19.61\% PCE
{[}\href{https://doi.org/10.1002/solr.202500906}{109}{]}.

Lattice microstrain from A-site mismatch divides stabilization
strategies between static coordination and dynamic mechanical
regulation. Static passivators such as guanidinium methanesulfonate
(GUAMS) bridge FA+ via hydrogen bonding while passivating \ce{Pb^2+},
achieving 26.11\% PCE in inverted cells, 24.49\% on flexible substrates,
and a T80 of \qty{1200}{\hour}
{[}\href{https://doi.org/10.1002/adfm.76197}{110}{]}. Zwitterionic
Cl-PPS additives identically coordinate \ce{Pb^2+} and hydrogen-bond
with FA+ to reach 26.61\% PCE
{[}\href{https://doi.org/10.1002/aenm.71105}{111}{]}. In n-i-p devices,
thiacetazone yields 26.52\% PCE in ambient air and 23.24\% at 70\% RH
{[}\href{https://doi.org/10.1002/adma.202523630}{112}{]}, while
fluoropyridine additives (DFPy) suppress non-radiative recombination to
reach 26\% PCE {[}\href{https://doi.org/10.1002/adfm.202526679}{113}{]}.
Because static bonds cannot accommodate cyclic thermal expansion,
dynamic strain regulation via thermal disulfide ring-opening crosslinks
grain boundaries, maintaining 95.7\% of initial efficiency after
1,\qty{800}{\hour} under maximum power point tracking (MPPT) with
day-night cycling (26.5\% champion PCE)
{[}\href{https://doi.org/10.1038/s41467-026-70697-5}{114}{]}.
Fluorinated misplaced-dipole elastomers similarly suppress thermal
stress cracks, yielding 25.54\% PCE on flexible devices versus 26.83\%
on rigid substrates
{[}\href{https://doi.org/10.1038/s41467-026-74285-5}{115}{]}. At module
scale, out-of-plane strain homogenization through A-site doping and an
in situ 2D stress-relief layer delivers a certified 25.55\% PCE (26.59\%
champion) across \qty{100}{\centi\meter\squared} under International
Summit on Organic Photovoltaic Stability (ISOS)-L-1 testing
{[}\href{https://doi.org/10.1126/sciadv.aec3238}{13}{]}. Solvent-free
dry vacuum deposition independently demonstrates that minimizing
microstrain yields a standard deviation of 0.36\% across 90 cells
(25.4\% champion PCE)
{[}\href{https://doi.org/10.1002/cey2.70236}{116}{]}.

In wide-bandgap (WBG) mixed-halide perovskites, light-induced halide
segregation divides mechanistic interpretations between differential
crystallization kinetics and lattice microstrain. Eutectic urea
phosphate demonstrates that halide demixing stems from unequal growth
kinetics; stronger affinity for bromine-based over iodine-based
octahedra synchronizes crystallization, eliminating vertical and lateral
phase segregation (22.27\% PCE)
{[}\href{https://doi.org/10.1002/anie.2535666}{117}{]}. TTA additives
corroborate this kinetic synchronization mechanism, achieving 20.55\%
PCE and an open-circuit voltage (Voc) of \qty{1.339}{\volt}
{[}\href{https://doi.org/10.1002/smll.202513111}{118}{]}. In contrast,
methylammonium-guanidinium (MA-GA) codoping demonstrates that balancing
lattice strain directly stabilizes the lattice against halide demixing
without kinetic agents, retaining over 95\% of initial efficiency after
\qty{1400}{\hour} of continuous operation in \qty{1.77}{\electronvolt}
cells (21.5\% PCE)
{[}\href{https://doi.org/10.1021/acsenergylett.6c00328}{119}{]}. Halide
substitution also dictates transport limits in thick absorbers: in
Cs0.05(FA0.90MA0.10)0.95Pb(I0.95Br0.05)3 middle cells, performance drops
above \qty{450}{\nano\meter}, but chloride substitution permits
thickness scaling up to \textasciitilde{}\qty{700}{\nano\meter} without
sacrificing fill factor (FF) or Voc (23.4\% PCE on
\qty{1}{\centi\meter\squared})
{[}\href{https://doi.org/10.1002/solr.70410}{120}{]}. Functional
group-modulated benzoic additives suppress non-radiative recombination
in WBG single junctions to exceed 22.9\% PCE, enabling 2-terminal tandem
cells surpassing 31\%
{[}\href{https://doi.org/10.1002/adfm.75640}{121}{]}. Multifunctional
ATAA molecules stabilize both FA+/MA+ cations and I-/Br- halides via
hydrogen bonding, upshifting the valence band maximum (VBM) by
\qty{0.12}{\electronvolt} (25.82\% PCE on \qty{1}{\centi\meter\squared})
{[}\href{https://doi.org/10.1002/adfm.202524302}{122}{]}, while dual
A-site 1D (PDA)0.6(GUA)1.4Pb6I16 bridges suppress iodide migration and
interfacial recombination, attaining 20.46\% PCE (Voc
\qty{1.328}{\volt}) in WBG cells and 28.43\% in all-perovskite tandems
with a T80 of \qty{1000}{\hour}
{[}\href{https://doi.org/10.1021/acsenergylett.5c04243}{123}{]}.

Low-dimensional heterojunctions introduce fundamental trade-offs between
surface passivation and transport resistance. Precursor incorporation of
long-chain oleylammonium ligands forms in situ 2D/3D heterojunctions
yielding 26.22\% PCE on small areas, 24.6\% on
\qty{1}{\centi\meter\squared}, and 21.1\% on
\qty{6.8}{\centi\meter\squared} mini-modules (T80 of \qty{1050}{\hour})
{[}\href{https://doi.org/10.1007/s40820-025-02058-8}{124}{]}.
Non-Dion-Jacobson (non-DJ) 2D perovskites using
3,3'-methylenediphenyldiammonium cations utilize organic bilayer pi-pi
interactions to achieve 26.52\% PCE on \qty{50}{\centi\meter\squared}
modules {[}\href{https://doi.org/10.1002/anie.2974414}{125}{]}. At
single-crystal interfaces, layered 1D (4PPY)PbI3 blocking layers
suppress iodide migration more effectively than isolated 1D or 2D
phases, retaining 100\% of initial efficiency after \qty{1032}{\hour}
continuous illumination (25.4\% PCE)
{[}\href{https://doi.org/10.1002/adfm.75386}{126}{]}. Alkyl chain
optimization with tetrahexylammonium iodide (THAI) generates 1D THAPbI3
heterojunctions that anchor FA+ via hydrogen bonding, suppressing
outward cation diffusion and strain gradients (25.06\% PCE)
{[}\href{https://doi.org/10.1021/acsaem.6c01363}{127}{]}. Grain-boundary
encapsulation via 4-FAPyCl2 in \ce{PbI2} precursors forms quasi-2D
coatings that protect 3D alpha-\ce{FAPbI3}, limiting ambient degradation
to less than 3\% PCE loss (25.32\% PCE)
{[}\href{https://doi.org/10.1002/adfm.202531852}{128}{]}. Precursor
addition of 0.3\% MeO-PEACl raises Voc from 1.11 to \qty{1.16}{\volt}
and PCE from 21.1\% to 23.1\%, maintaining the black phase after
\qty{3380}{\hour} at 65\% ± 10\% RH
{[}\href{https://doi.org/10.1002/smll.73351}{129}{]}. To overcome spacer
insulation, thiophene-2-hydrazide (ThCH) induces interlayer orbital
coupling in 2D Ruddlesden-Popper (RP) phases, though its certified
21.74\% efficiency (22.41\% champion) was evaluated on a diminutive
\qty{0.074}{\centi\meter\squared} aperture
{[}\href{https://doi.org/10.1002/anie.4293157}{130}{]}. Pre-formed n=1
2D nanosheets seed 3D oriented growth to reach 21.36\% PCE (20.75\%
stabilized output) with hysteresis below 1\%
{[}\href{https://doi.org/10.1002/smll.74076}{131}{]}, and
poly(3-hexylthiophene) (P3HT)-modified 2D quantum dots suppress
interfacial defects to reach 25.04\% PCE
{[}\href{https://doi.org/10.1021/acsami.5c24915}{132}{]}. Isomeric
cytosine additives outperform iso-cytosine through delocalized bonding
in quasi-2D films (22.4\% PCE; T80 of \qty{3600}{\hour})
{[}\href{https://doi.org/10.1021/acs.jpclett.6c01531}{133}{]}, and
propylammonium iodide in two-step blade coating directs (111) grain
orientation in inverted cells (20.05\% PCE; 95\% retention after
\qty{1680}{\hour} MPPT)
{[}\href{https://doi.org/10.1002/solr.70385}{134}{]}. Conversely,
unoptimized capping layers such as 2D (4-FBA)2PbBr4 yield an uncertified
18.4\% PCE and retain only 77.82\% after \qty{10}{\hour} at 80 °C and
55.32\% after 30 days in ambient air
{[}\href{https://doi.org/10.1021/acsaem.6c00285}{135}{]}, while vacuum
co-evaporated PEA-passivated films reach 17.79\% PCE
{[}\href{https://doi.org/10.1002/ente.202502217}{136}{]}.

\subsection{6. Tandem Architecture and Current
Matching}\label{tandem-architecture-and-current-matching}

Investigations into monolithic two-terminal (2T) perovskite/silicon
tandems diverge on whether a laterally conductive transparent conducting
oxide (TCO) interlayer is required for subcell recombination. Conformal
coverage on textured silicon heterojunction (SHJ) bottom cells
reportedly required an \qty{8}{\nano\meter} indium tin oxide (ITO) layer
modified with nickel oxide (\ce{NiO_x}) and a polymeric self-assembled
monolayer (Poly-SAM), gaining
\qty{0.85}{\milli\ampere\per\centi\meter\squared} in short-circuit
current density (JSC) to reach
\qty{20.82}{\milli\ampere\per\centi\meter\squared} in a 31.8\% champion
cell {[}\href{https://doi.org/10.1002/adfm.75311}{137}{]}. Conversely,
omitting indium-based oxides entirely, a reactive plasma deposition
(RPD) tin oxide (SnOx) recombination layer directly anchored
self-assembled monolayers (SAMs) to suppress halide migration and
nonradiative recombination, achieving certified 33.6\% power conversion
efficiency (PCE) (33.2\% champion on \qty{1}{\centi\meter\squared}) and
certified 31.0\% on a \qty{207.9}{\centi\meter\squared} minimodule
{[}\href{https://doi.org/10.1126/science.aef5355}{1}{]}. The physical
controversy centers on whether removing ITO lateral conductivity creates
localized resistive bottlenecks across textured pyramids. Spatially
resolved voltage-dependent photoluminescence (PL) imaging via the local
transport and recombination study (LTRS) method settles this debate: by
mapping series resistance (Rs) to isolate an absolute 1.9\% efficiency
loss in a 29\% monolithic tandem while decoupling current mismatch, low
shunt resistance, and luminescence coupling in finite element
simulations, operando LTRS determines whether TCO-free junctions
exacerbate resistive losses under spectral variation
{[}\href{https://doi.org/10.1002/smll.202513958}{138}{]}.

Parasitic optical absorption within selective contacts further
constrains current matching. Thinning electron-selective fullerene
(\ce{C60}) contacts to \qty{1}{\nano\meter} yields thickness-independent
extraction, boosting tandem JSC by
\qty{0.24}{\milli\ampere\per\centi\meter\squared} in a 31.7\% champion
cell {[}\href{https://doi.org/10.1021/acsenergylett.5c04082}{139}{]}.
Yet atomic layer deposition (ALD) of tin oxide damages underlying
perovskite during \ce{C60} capping; substituting ALD with thermally
evaporated antimony oxide eliminates damage while permitting thinner
fullerenes, improving top-cell JSC by
\textasciitilde{}\qty{1}{\milli\ampere\per\centi\meter\squared}
(certified 27.7\%, champion 28.16\% on
\qty{64.64}{\centi\meter\squared})
{[}\href{https://doi.org/10.1038/s41467-026-70848-8}{140}{]}. In normal
(n-i-p) tandems, an ultrathin \qty{10}{\nano\meter} dopant-free
2,2',7,7'-tetrakis{[}N,N-di(4-methoxyphenyl)amino{]}-9,9'-spirobifluorene
(spiro-OMeTAD) layer reduces parasitic absorption by 92.2\% and
reflection losses by 18.4\%, delivering an uncertified 29.73\% PCE on
\qty{0.135}{\centi\meter\squared} and 28.77\% (certified 28.25\%) on
\qty{1.012}{\centi\meter\squared}
{[}\href{https://doi.org/10.1007/s40820-025-01962-3}{141}{]}. Decoupling
current matching, four-terminal (4T) architectures using semitransparent
\qty{1.67}{\electronvolt} wide-bandgap (WBG) subcells with vanadium
oxide-nickel oxide (VOx-\ce{NiO_x}) contacts achieved 21.68\% on
\qty{0.148}{\centi\meter\squared}, retaining 98\% efficiency after
\qty{1472}{\hour} maximum power point tracking (MPPT) at 25 °C (31.25\%
tandem on \qty{1}{\centi\meter\squared})
{[}\href{https://doi.org/10.1002/aenm.71036}{142}{]}, while priming
CsPbX3 intermediate phases yielded 20.91\% on
\qty{0.1}{\centi\meter\squared} and 19.67\% on
\qty{1.0}{\centi\meter\squared} (31.04\% tandem on
\qty{1.0}{\centi\meter\squared})
{[}\href{https://doi.org/10.1038/s41467-026-72094-4}{143}{]}.

Matching bottom cells requires WBG absorbers
(1.67-\qty{1.85}{\electronvolt}), where drift-diffusion modeling
indicates that bromide substitution and cesium alloying in formamidinium
lead triiodide (\ce{FAPbI3}) increase carrier effective mass from 0.173
me to 0.243-0.258 me, restricting transport despite projected 32.10\%
PCE and \qty{2.05}{\volt} open-circuit voltage (VOC)
{[}\href{https://doi.org/10.26434/chemrxiv.15001694/v1}{144}{]}.
Approaches over the six-month period diverge on additive chemistry.
Several groups introduced fluorinated agents: difluoro(oxalato)borate
and tetrafluoroborate synchronized crystallization without intermediate
halide redistribution, certifying a 30.3\% tandem
{[}\href{https://doi.org/10.1038/s41565-026-02165-6}{145}{]}; volatile
octafluoronaphthalene (OFNA) relieved tensile stress and passivated
defects (22.23\% single junction, 31.43\% tandem)
{[}\href{https://doi.org/10.1002/adfm.202526009}{146}{]}; and
4-(trifluoromethyl)benzenesulfonyl chloride (BSCF3) coordinated lead and
formamidinium (19.81\% single junction, \qty{1.32}{\volt} VOC, 27.26\%
tandem) {[}\href{https://doi.org/10.1002/solr.70387}{147}{]}. In direct
conflict, tracking ion diffusion revealed that mobile fluoride from
lithium fluoride (LiF) interlayers penetrates the perovskite bulk,
accelerating halide segregation and degrading photostability; an
ion-blocking lithium-theophylline (Li-T) interlayer was required to halt
fluoride migration, securing \qty{1.284}{\volt} VOC and 19.55\%
efficiency in \qty{1.77}{\electronvolt} cells (28.41\% tandem)
{[}\href{https://doi.org/10.1002/aenm.202506634}{148}{]}.
Non-fluorinated strategies instead employed forchlorfenuron (CPPU) to
direct (100) facets (certified 23.46\% single junction,
\qty{1.295}{\volt} VOC; certified 31.93\%, champion 32.17\% tandem on
\qty{1.0}{\centi\meter\squared})
{[}\href{https://doi.org/10.1016/j.esci.2026.100563}{3}{]},
1,2-diiodotetrafluorobenzene (1,2-DITFB) to retard bromide
crystallization (20.25\% single junction, \qty{1.34}{\volt} VOC, 28.58\%
tandem) {[}\href{https://doi.org/10.1002/adfm.74659}{149}{]},
phenylhydrazine derivatives to inhibit iodide oxidation (20.75\% single
junction, 28.63\% tandem)
{[}\href{https://doi.org/10.1002/smll.202510652}{150}{]}, and
2-phenylethylammonium chloride (PEACl) modification with solvent etching
(CDSE) to remove residual lead iodide (\ce{PbI2}) (21.77\% single
junction, 31.85\% tandem)
{[}\href{https://doi.org/10.1002/aenm.202506619}{151}{]}. Rubidium
cyanate (RbOCN) lattice incorporation stabilized
\qty{1.73}{\electronvolt} absorbers (22.45\% single junction, 25.37\% on
\qty{1}{\centi\meter\squared}, 26.75\% champion on an uncertified
\qty{0.09}{\centi\meter\squared} area, \qty{2.22}{\volt} VOC)
{[}\href{https://doi.org/10.1002/adma.72986}{152}{]}, while targeted
passivator coordination achieved a certified \qty{2.00}{\volt} tandem
VOC (certified 31.52\%, champion 32.74\%)
{[}\href{https://doi.org/10.1002/aenm.71077}{153}{]}. In contrast,
additive-free stoichiometry control via a \ce{PbI2}-to-total PbX2 ratio
of \textasciitilde55-75\% slowed crystallization to deliver
single-junction PCEs of \textasciitilde24.19\%
(\qty{1.67}{\electronvolt}) to 16.83\% (\qty{1.93}{\electronvolt}) and a
29.1\% tandem {[}\href{https://doi.org/10.1021/acsnano.6c04740}{154}{]}.
Evaporated scaffolds expose a mechanistic disconnect: two-step thermal
evaporation yielded stabilized tandem PCEs above 29\%
{[}\href{https://doi.org/10.1021/acsenergylett.6c00156}{155}{]}, yet
five sequentially evaporated scaffold stacks produced identical
\textasciitilde19\% single-junction PCEs (28.4\% champion tandem)
despite profound variations in precursor conversion and morphology
{[}\href{https://doi.org/10.1021/acsami.6c07582}{156}{]}.

At buried hole-selective junctions, composite SAMs of
4,4',4'\,`,4'\,'\,'-methanetetrayltetrabenzoic acid (4MA) with
{[}4-(3,6-dimethyl-9H-carbazol-9-yl)butyl{]}phosphonic acid (Me-4PACz)
on \ce{NiO_x} use tetrahedral molecular pinning to passivate defects,
achieving 23.52\% in \qty{1.68}{\electronvolt} cells and a certified
32.1\% (32.26\% champion on \qty{1}{\centi\meter\squared}) tandem
{[}\href{https://doi.org/10.1002/anie.5162616}{12}{]}. Aligning packing
via intermolecular pi-pi interactions in
{[}2-(3,6-dimethyl-9H-carbazol-9-yl)ethyl{]}diphenylphosphine oxide
(Me-Ph2mPACz) reduced non-radiative losses from \qty{168}{\milli\volt}
to \qty{124}{\milli\volt} (single-junction PCE 21.82\% to 23.14\%),
delivering a certified 32.45\% (33.4\% champion) tandem retaining 83\%
efficiency after \qty{1000}{\hour} MPPT at 85 °C in air
{[}\href{https://doi.org/10.1038/s41467-026-72160-x}{8}{]}.
Three-dimensional C-H\ldots pi transport networks from asymmetric
MeO-BADCz SAMs certified at 33.04\% (24.05\% champion reported)
{[}\href{https://doi.org/10.1038/s41467-026-73656-2}{11}{]}, and the
triazole-based quaternary ammonium molecule TMTEA co-regulated
self-assembled monolayer (SAM) packing and orientation (23.93\% single
junction, 32.15\% champion)
{[}\href{https://doi.org/10.1021/acsenergylett.6c00166}{157}{]}. Under
operational current mismatch, subcells are driven into reverse bias. In
all-perovskite tandems combining WBG and narrow-bandgap (NBG) absorbers,
sustained reverse bias induces a destructive redox cascade: hole
injection oxidizes iodide to iodine species that oxidize \ce{Sn^2+} to
\ce{Sn^4+}, driving lattice collapse in a 29.03\% tandem
{[}\href{https://doi.org/10.1002/adma.202522942}{158}{]}.

\subsection{7. Degradation Mechanisms Under
Operation}\label{degradation-mechanisms-under-operation}

Studies of operational degradation reveal an unresolved mechanistic
contradiction regarding whether bulky molecular passivators relieve
interfacial lattice strain or accelerate degradation through local
thermal expansion. Bulky passivators are widely applied to terminate
dangling bonds and impede defect formation, with molecular locks such as
3-PMPCl raising the phase transition barrier in formamidinium lead
triiodide (\ce{FAPbI3}) by suppressing {[}PbI6{]} octahedral disorder to
reach a certified 27.6\% power conversion efficiency (PCE)
{[}\href{https://doi.org/10.1126/science.aeb9953}{159}{]}. Similarly,
facet-engineering additives such as sulfonates relieve residual lattice
strain by redirecting crystallization kinetics
{[}\href{https://doi.org/10.1002/adma.73539}{160}{]}, and molecular
buffering layers prevent mechanochemical fatigue under thermal cycling
{[}\href{https://doi.org/10.1002/adma.202519339}{161}{]}. However,
direct operando structural analysis demonstrates that bulky anionic
passivators induce local lattice expansion under thermal stress at
heterointerfaces, driving interfacial fatigue even while unencapsulated
devices retain 87\% of their initial efficiency after \qty{4092}{\hour}
of maximum power point tracking (MPPT) under International Summit on
Organic Photovoltaic Stability (ISOS) protocol ISOS-L-1 with a time to
80\% of initial performance (T80) of \qty{3566}{\hour}
{[}\href{https://doi.org/10.1126/sciadv.aef6596}{162}{]}. This finding
directly challenges the assumption that bulky organic modifiers serve as
inert physical anchors. The contradiction cannot be resolved by
post-mortem spectroscopy because lattice relaxation occurs upon cooling;
settling it requires in situ, temperature-dependent grazing-incidence
wide-angle X-ray scattering paired with depth-resolved photoelectron
spectroscopy under operational bias to quantify local lattice expansion
directly against defect formation rates.

A parallel divergence divides studies of buried hole-selective
interfaces in inverted (p-i-n) architectures. One body of literature
identifies out-of-plane molecular desorption from metal oxides as the
primary degradation mode, which can be mitigated by enforcing spatial
confinement with a rigid planar backbone achieving 3.72 Å pi-pi stacking
{[}\href{https://doi.org/10.1016/j.scib.2026.02.050}{163}{]} or by
synthesizing polymeric self-assembled monolayers (SAMs) such as
Poly-4PADCB to reduce steric hindrance
{[}\href{https://doi.org/10.1002/adma.72622}{164}{]}. Conversely,
competing evidence attributes contact failure to spontaneous molecular
aggregation that leaves bare oxide domains, an issue countered through
asymmetric molecules with steric hindrance like PhBr-4PACz that retain
over 96\% efficiency after \qty{2000}{\hour} at 65 °C under ISOS-L-2
{[}\href{https://doi.org/10.1038/s41467-026-72097-1}{165}{]} or
fluorinated confinement interlayers
{[}\href{https://doi.org/10.1038/s41467-026-72115-2}{166}{]}. A third
perspective indicates that self-assembled monolayer (SAM) degradation
occurs predominantly during perovskite solution processing via solvent
stripping, requiring transient 2D perovskite sacrificial barriers
{[}\href{https://doi.org/10.1002/adma.73561}{167}{]} or hydroxyl
enrichment using MXene additives
{[}\href{https://doi.org/10.1002/adma.73656}{168}{]}. Although amino
acid hydrochloride additives
{[}\href{https://doi.org/10.1002/aenm.70925}{169}{]} and tailored hole
transport materials like WH13
{[}\href{https://doi.org/10.1002/anie.202523665}{170}{]} yield stability
claims, these evaluations frequently rely on tiny active areas (such as
\qty{0.058}{\centi\meter\squared} in
{[}\href{https://doi.org/10.1002/aenm.70925}{169}{]}), where edge
encapsulation masks real-world planar degradation. In perovskite/silicon
tandem cells, co-assembled siloxane monolayers
{[}\href{https://doi.org/10.1002/aenm.71107}{10}{]} and amorphous
asymmetric contacts {[}\href{https://doi.org/10.1002/adma.73106}{6}{]}
retain 87.12\% and 87\% initial PCE after \qty{1000}{\hour} of MPPT at
85 °C, respectively, confirming that resilient chemisorption retards
thermally activated interfacial delamination.

In the perovskite absorber, the consensus that device decay is governed
primarily by bulk mobile ions has been contested by direct
quantification of deep-level defect kinetics. Operational degradation is
dominated by deep-level interstitial and antisite defects, specifically
iodine on formamidinium sites (I\_FA) and lead sites (I\_Pb), which
drive performance decay despite existing at concentrations three orders
of magnitude below shallow defects
{[}\href{https://doi.org/10.1002/adma.73679}{171}{]}. Passivating these
deep traps with 3TU2+ reduces rear-interface energy losses from 1.46\%
to 0.62\% {[}\href{https://doi.org/10.1002/adma.73679}{171}{]}. In
contrast, widespread reports continue to focus on macroscopic ionic
flux, using liquid additives such as 1-BP to suppress ion migration as
tracked by time-of-flight secondary ion mass spectrometry (ToF-SIMS)
{[}\href{https://doi.org/10.1038/s41467-025-68231-0}{172}{]}, positional
methylation isomers like 5-MH that retain 92\% efficiency after
\qty{2000}{\hour} under ISOS-L2
{[}\href{https://doi.org/10.1002/adma.202520432}{173}{]}, or eperisone
hydrochloride to eliminate residual dimethyl sulfoxide (DMSO) and
maintain 90.39\% efficiency after \qty{1200}{\hour} under ISOS-L-3
{[}\href{https://doi.org/10.1021/acsenergylett.5c03999}{174}{]}. Other
interventions combat defect strain through trifluoroethylammonium
chloride (3TFACl) additives
{[}\href{https://doi.org/10.1002/adfm.76756}{175}{]} or 1D (PCOZ)PbI3
interlayers that preserve 94.2\% efficiency after \qty{1000}{\hour}
under ISOS-L-3 {[}\href{https://doi.org/10.1002/adma.73161}{176}{]}.

Operational stability is also restricted by charge transport layers,
where conflicting failure mechanisms are assigned to identical
materials. In inverted devices, phenyl-C61-butyric acid methyl ester
(PCBM) degrades via photo-induced {[}2+2{]} cycloaddition dimerization
under illumination, inducing aggregation and morphological decay that
nonplanar tetraphenylethylene additives prevent
{[}\href{https://doi.org/10.1002/solr.70405}{177}{]}. Yet separate
studies treat fullerene decline as an interfacial corrosion problem
solved by buffer layers like BAE-Phen that coordinate migrant metal ions
to retain 90.5\% PCE after \qty{2000}{\hour} at 85 °C
{[}\href{https://doi.org/10.1002/adma.73529}{178}{]}, solvent-free
multicomponent fullerene composites
{[}\href{https://doi.org/10.1038/s41467-026-70022-0}{179}{]}, or
bridging interlayers tested on \qty{0.045}{\centi\meter\squared} cells
with a T80 of \qty{500}{\hour}
{[}\href{https://doi.org/10.1002/anie.9286497}{180}{]}. In normal
(n-i-p) devices employing
2,2',7,7'-tetrakis{[}N,N-di(4-methoxyphenyl)amino{]}-9,9'-spirobifluorene
(Spiro-OMeTAD), a fundamental trade-off persists: promoting lithium
bis(trifluoromethanesulfonyl)imide (LiTFSI) ion-pair dissociation is
necessary to sustain electrical conductivity
{[}\href{https://doi.org/10.1002/aenm.202506760}{181}{]}, yet
uncoordinated lithium ions migrate rapidly into the absorber.
Suppressing this degradation channel requires either phenothiazine ionic
anchors to bind lithium, yielding 90\% retention after \qty{800}{\hour}
of MPPT on \qty{0.04}{\centi\meter\squared} cells
{[}\href{https://doi.org/10.1002/adfm.76257}{182}{]}, or lithium-based
deep eutectic matrices that deliver 90.5\% retention after
\qty{1800}{\hour} under ISOS-L-1I
{[}\href{https://doi.org/10.1002/anie.9551164}{183}{]}.

These mechanistic debates unfold against an acute reporting deficit
across the field: of 872 closely read papers published over the half
year, only 39 reported a T80 value and only 39 specified an ISOS testing
protocol. Although champion operational lifespans reached
\qty{4100}{\hour} {[}\href{https://doi.org/10.1002/adfm.76110}{184}{]},
\qty{4000}{\hour} {[}\href{https://doi.org/10.1002/adfm.74893}{185}{]},
\qty{3600}{\hour}
{[}\href{https://doi.org/10.1021/acs.jpclett.6c01531}{133}{]}, and
\qty{3566}{\hour}
{[}\href{https://doi.org/10.1126/sciadv.aef6596}{162}{]}, the rarity of
standardized aging procedures leaves the relative severity of
interfacial expansion, ion accumulation, and transport layer degradation
obscured by divergent testing metrics.

\begin{figure}[htbp]
\centering
\includegraphics[width=0.95\textwidth]{fig/F4_stability_evidence.pdf}
\caption{The month's operational-stability evidence in full. Left, every reported time to 80\% of initial performance. Right, the stability protocol labels used.}
\end{figure}

\subsection{8. Large-Area Coating and the Area
Penalty}\label{large-area-coating-and-the-area-penalty}

Conflicting kinetic models across recent scalable coating studies
reflect an unresolved physical disagreement over whether
precursor-solvent coordination accelerates or impedes uniform perovskite
crystallization. One framework asserts that prolonged solvent-precursor
interaction time during blade coating stabilizes detrimental
solvent-adduct intermediate phases that degrade film crystallinity,
demanding rapid adduct suppression
{[}\href{https://doi.org/10.1038/s41467-026-68439-8}{186}{]}. In direct
contradiction, alternative investigations conclude that strong
coordination by N-methyl-2-pyrrolidone (NMP) in dimethylformamide
(DMF)-free 2-methoxyethanol (2-ME) or isopropanol (IPA) is necessary to
stabilize intermediate complexes and delay premature surface nucleation,
thereby increasing open-circuit voltage (Voc) from 1.08 to
\qty{1.09}{\volt} and fill factor (FF) from 75\% to 80.5\%
{[}\href{https://doi.org/10.1063/5.0297564}{187},\href{https://doi.org/10.1002/solr.70406}{188}{]}.
This kinetic divergence extends across solvent systems: ternary dimethyl
sulfoxide (DMSO), acetonitrile (ACN), and 1,3-dimethyl-2-imidazolidinone
(DMI) mixtures rely on volatile ACN to accelerate initial nucleation
during slot-die printing to reach 20.62\% power conversion efficiency
(PCE) {[}\href{https://doi.org/10.1016/j.matre.2026.100406}{189}{]},
whereas reduced-temperature bar coating at 100 °C employs low
translation speeds of 0.3 mm/s to sustain a metastable supersaturation
zone rather than accelerating precipitation
{[}\href{https://doi.org/10.1002/solr.202600003}{190}{]}. Resolving
whether solvent adducts degrade or preserve large-area films requires
operando spatial diagnostics; coupling in situ photoluminescence (PL)
and diffuse reflectance during quenching with neural networks captures
non-uniform crystallization before localized shunts solidify
{[}\href{https://doi.org/10.3390/mi17030314}{191},\href{https://doi.org/10.1007/s40820-026-02134-7}{192}{]}.
Precursor colloid engineering via formate anions restricts cluster size
to around \qty{30}{\nano\meter}
{[}\href{https://doi.org/10.1002/adma.202520776}{193}{]}, and
polypropylene glycol (PPG) aids formamidinium lead iodide (\ce{FAPbI3})
microcrystal synthesis
{[}\href{https://doi.org/10.1002/adma.202522508}{194}{]}. However,
claims of defect suppression through adduct control remain restricted to
uncertified, sub-\qty{0.1}{\centi\meter\squared} cells (such as 26.5\%
PCE on \qty{0.0665}{\centi\meter\squared}
{[}\href{https://doi.org/10.1038/s41467-026-68439-8}{186}{]} or
certified 26.22\% on \qty{0.06}{\centi\meter\squared}
{[}\href{https://doi.org/10.1002/adma.202522508}{194}{]}), without
proving that intermediate coordination equilibria survive roll-to-roll
drying gradients under mechanical shear
{[}\href{https://doi.org/10.20517/ss.2025.128}{195}{]}.

A second major source of the area penalty originates in charge transport
layer (CTL) non-uniformity and buried contact resistance across large
apertures. For electron transport layers (ETLs), atmospheric-pressure
spatial atomic layer deposition (AP-SALD) of tin oxide (SnOx) achieves
20\% PCE following 180 °C annealing, yet conformal deposition replicates
substrate roughness, which impairs perovskite crystallization
{[}\href{https://doi.org/10.1016/j.solmat.2026.114233}{196}{]}. Solution
approaches attempt to mitigate interfacial roughness: electrostatic
regulation of \ce{SnO2} nanoparticle assembly yields uniform coatings,
delivering 24.70\% PCE at \qty{1}{\centi\meter\squared} and over 21\% on
a \qty{65.24}{\centi\meter\squared} module (certified 26.03\% on
\qty{0.0727}{\centi\meter\squared})
{[}\href{https://doi.org/10.1021/acsenergylett.6c00476}{197}{]}, while
multicycle infrared annealing of \ce{SnO2} quantum dot (QD) ETLs limits
cracking, reaching 17.09\% at \qty{1}{\centi\meter\squared} and 15.18\%
at \qty{2}{\centi\meter\squared} compared to 12.78\% for conventional
hotplate annealing
{[}\href{https://doi.org/10.1016/j.rineng.2026.111690}{198}{]}. For
hole-selective contacts, self-assembled monolayers (SAMs) suffer severe
dewetting during large-area coating; incorporating 3-BPIC-F into
MeO-4PACz enhances dipole alignment and suppresses pinholes, raising
module PCE from 16.43\% to 19.80\% (champion 23.64\%)
{[}\href{https://doi.org/10.1021/acsami.6c03670}{199}{]}, while
low-polarity alkane and alcohol solvents enable blade coating of SAMs at
80\% relative humidity (RH) to yield a certified 26.1\% PCE
{[}\href{https://doi.org/10.1038/s41467-026-72581-8}{200}{]}. At the
counter interface, non-uniform chemical doping induces spatial
conductivity variances, which ferrocenium radical cation dopants in
Spiro-OMeTAD eliminate to achieve 26.13\% in cells and 22.21\% in
modules {[}\href{https://doi.org/10.1016/j.joule.2025.102310}{201}{]}.

Hydrodynamic instability and resistive interconnect losses govern the
final translation from single devices to large modules. Drying flows
induce severe thickness gradients; fluid-locking interfacial layers
counteract coffee-ring segregation to yield 21.39\% PCE on a
\qty{100}{\centi\meter\squared} module
{[}\href{https://doi.org/10.1002/adma.202518582}{202}{]}, whereas
slot-die coated \ce{FAPbI3} without fluid control drops from 20.8\% on
small areas to 19.1\% at \qty{8.4}{\centi\meter\squared} and 14.8\% at
\qty{59.5}{\centi\meter\squared} active area
{[}\href{https://doi.org/10.1021/acsami.5c22032}{203}{]}.
Post-deposition passivation attempts to seal boundary defects: immersing
hot-bladed films into cold strontium iodide (SrI2) isopropanol solutions
produces 20.32\% on a \qty{10.13}{\centi\meter\squared} module
{[}\href{https://doi.org/10.1038/s41467-026-69264-9}{204}{]}, and
pneumatic spray coating forms 1D PyPbI3 surface barriers (champion
21.64\%) {[}\href{https://doi.org/10.1016/j.mtsust.2026.101347}{205}{]}.
Fully dry vacuum co-evaporation avoids wet-film fluid mechanics
entirely, using \ce{PbCl2} additions to disrupt \ce{PbI2} stacking and
reach over 38\% efficiency at 900 lux on a
\qty{3.9}{\centi\meter\squared} indoor module
{[}\href{https://doi.org/10.1002/smtd.70654}{206}{]}. Beyond
\qty{100}{\centi\meter\squared}, mechanical deformation and interconnect
resistance dominate: dynamic covalent cross-linking monomers embedded in
perovskites provide self-curing at 40 °C for \qty{30}{\minute}, yielding
a certified 26.8\% and champion 27.12\% PCE across
\qty{655.2}{\centi\meter\squared} flexible modules
{[}\href{https://doi.org/10.1002/adma.73855}{15}{]}, while liquid metal
interconnects maintain continuity over a \qty{900}{\centi\meter\squared}
platform {[}\href{https://doi.org/10.1002/adfm.202532205}{18}{]}.

\begin{figure}[htbp]
\centering
\includegraphics[width=0.95\textwidth]{fig/F2_area_penalty.pdf}
\caption{Efficiency against reported active or aperture area. The shaded band marks the sub-0.3 cm$^2$ laboratory-cell regime in which most champion values are measured.}
\end{figure}

\subsection{9. Cross-Cutting Synthesis: What the Mechanisms Say
Together}\label{cross-cutting-synthesis-what-the-mechanisms-say-together}

A fundamental disagreement across recent multi-junction investigations
concerns whether non-radiative voltage deficits originate from defective
charge extraction at buried contacts or from trap-assisted recombination
within bulk wide-bandgap (WBG) absorbers. In monolithic two-terminal
(2T) perovskite/silicon tandems, one framework attributes performance
bottlenecks primarily to hole-extraction resistance and halide
accumulation at the buried contact, demonstrating that replacing indium
tin oxide (ITO) with a reactive plasma deposition (RPD) tin oxide (SnOx)
recombination layer directly anchored to self-assembled monolayers
(SAMs) delivers a certified power conversion efficiency (PCE) of 33.6\%
(champion 33.2\% on \qty{1}{\centi\meter\squared}) and a certified
31.0\% on a \qty{207.9}{\centi\meter\squared} minimodule
{[}\href{https://doi.org/10.1126/science.aef5355}{1}{]}. Conversely,
competing investigations conclude that carrier losses are dominated by
photoinduced phase segregation and deep-level defect generation in the
bulk WBG film, asserting that replacing monomeric molecules with
polymeric polycarbazole phosphonic acid (Poly-3Ph-4PACz) relieves
bulk-driven recombination to reach a certified 2T tandem PCE of 32.95\%
(champion 26.13\%, stabilised 27.12\% in single junctions; 22.22\% on a
\qty{57.6}{\centi\meter\squared} flexible module)
{[}\href{https://doi.org/10.1038/s41467-026-73276-w}{2}{]}. Standard
steady-state photoluminescence (PL) and current density-voltage
characterization cannot decouple these hypotheses because contact
recombination velocity and bulk quasi-Fermi level splitting collapse
produce identical terminal open-circuit voltage (Voc) losses. Deciding
whether the primary voltage deficit resides at the contact or in the
absorber requires operando, spatially resolved voltage-dependent
photoluminescence imaging via the local transport and recombination
study (LTRS) method coupled with bias-dependent transient absorption
spectroscopy. By mapping series resistance to isolate a 1.9\% efficiency
loss in a 29\% monolithic tandem while decoupling current mismatch, low
shunt resistance, and luminescence coupling in finite-element
simulations, operando LTRS determines whether carrier losses originate
from interface extraction barriers or bulk non-radiative recombination
{[}\href{https://doi.org/10.1002/smll.202513958}{138}{]}.

Reading the separate axes together reveals that performance gains
claimed in one domain routinely impose physical penalties in another. At
the perovskite top surface, constructing low-dimensional capping
overlayers or deploying bulky organic ammonium passivators terminates
dangling bonds and heals halide vacancies, achieving certified PCE
values of 26.2\% (champion 26.49\%) with a
two-dimensional/three-dimensional (2D/3D) guanidinium heterostructure
{[}\href{https://doi.org/10.1002/anie.8890696}{20}{]}, 26.15\% (champion
26.71\% on \qty{1}{\centi\meter\squared}) through one-dimensional (1D)
TEAPbI3 phase reconstruction
{[}\href{https://doi.org/10.1038/s41467-026-72790-1}{19}{]}, and a
certified 27.6\% via 3-PMPCl molecular locking
{[}\href{https://doi.org/10.1126/science.aeb9953}{159}{]}. However,
evaluating these capping layers against operational degradation
demonstrates an incompatible physical demand: bulky organic cations
introduce anisotropic out-of-plane transport barriers that trap
photogenerated carriers under operating electric fields, while operando
structural measurements reveal that bulky passivators undergo pronounced
local thermal expansion relative to the perovskite lattice. Under
thermal stress, this expansion generates localized microstrain and
interfacial fatigue, driving degradation even in unencapsulated devices
that retain 87\% of initial efficiency after \qty{4092}{\hour} of
maximum power point tracking (MPPT) under International Summit on
Organic Photovoltaic Stability (ISOS) protocol ISOS-L-1 with a time to
80\% of initial performance (T80) of \qty{3566}{\hour}
{[}\href{https://doi.org/10.1126/sciadv.aef6596}{162}{]}. Incorporating
facet-engineering additives
{[}\href{https://doi.org/10.1002/adma.73539}{160}{]} or molecular
buffering interlayers
{[}\href{https://doi.org/10.1002/adma.202519339}{161}{]} to relieve
residual lattice strain demonstrates that chemical defect passivation
cannot be decoupled from thermomechanical fatigue: chemical saturation
of surface coordination bonds trades immediate electronic gains for
latent operational instability under thermal cycling.

An analogous conflict governs the buried interface, where electrostatic
optimization clashes with mechanical adhesion and lateral conductivity.
Functionalizing carbazole phosphonic acid SAMs with polar cyanoacetic
acid co-adsorbents nearly doubles the reverse-bias breakdown voltage to
yield an uncertified tandem champion PCE of 32.43\% (stabilised 24.43\%)
{[}\href{https://doi.org/10.1002/adma.72738}{87}{]}, while asymmetric
Br-4PhCz monolayers deliver a certified tandem PCE of 28.98\% (champion
29.35\%) on \qty{0.05}{\centi\meter\squared}
{[}\href{https://doi.org/10.1002/adma.202523338}{63}{]}. Yet monomeric
SAMs exhibit low fracture resistance and undergo chemical oxidation and
mechanical delamination under operational stress. Polymeric crosslinking
mitigates interfacial shear failure, as multidentate polymer networks
(P-SAM) increase interfacial fracture toughness nearly eightfold to 9.11
MPa (certified 26.61\% on \qty{0.06734}{\centi\meter\squared})
{[}\href{https://doi.org/10.1002/aenm.71159}{80}{]} and
phosphonic-acid-functionalized polymers (PTPP) sustain 94.4\% of initial
efficiency after 800 thermal cycles (certified 26.4\%, champion 27.1\%)
{[}\href{https://doi.org/10.1002/adfm.75697}{81}{]}. Nevertheless,
crosslinked polymeric interlayers create lateral transport resistance
across textured silicon bottom cells, where conformal coverage on random
pyramids requires either an \qty{8}{\nano\meter} ITO interlayer with
nickel oxide and polymeric SAMs (gaining
\qty{0.85}{\milli\ampere\per\centi\meter\squared} in short-circuit
current density to reach
\qty{20.82}{\milli\ampere\per\centi\meter\squared} in a 31.8\% champion
tandem) {[}\href{https://doi.org/10.1002/adfm.75311}{137}{]} or direct
self-assembled monolayer (SAM) anchoring on SnOx
{[}\href{https://doi.org/10.1126/science.aef5355}{1}{]}. Contact design
therefore cannot optimize dipole energetics, mechanical adhesion, and
lateral charge extraction simultaneously.

Finally, precursor coordination kinetics reflect an unbridged disconnect
between sub-\qty{0.1}{\centi\meter\squared} efficiency records and
scalable manufacturing. While non-destructive contact-triggered cationic
interaction (CCI) yields a certified 25.61\% PCE (champion 26.25\%)
{[}\href{https://doi.org/10.1038/s41560-026-02027-4}{98}{]} and
microcrystal growth with polypropylene glycol achieves a certified
26.22\% on \qty{0.06}{\centi\meter\squared}
{[}\href{https://doi.org/10.1002/adma.202522508}{194}{]},
solvent-coordination models reach diametrically opposing conclusions.
One framework demands rapid suppression of intermediate solvent adducts
to avoid structural defects (uncertified 26.5\% on
\qty{0.0665}{\centi\meter\squared})
{[}\href{https://doi.org/10.1038/s41467-026-68439-8}{186}{]}, whereas
another insists that strong coordination by N-methyl-2-pyrrolidone is
required to delay premature nucleation, increasing Voc from 1.08 to
\qty{1.09}{\volt} and fill factor (FF) from 75\% to 80.5\%
{[}\href{https://doi.org/10.1063/5.0297564}{187}{]}. These chemical
mechanisms do not converge because the intermediate equilibria that
suppress defect formation in static, small-area spin coating fail under
the mechanical shear and drying gradients of scalable deposition
{[}\href{https://doi.org/10.20517/ss.2025.128}{195}{]}. Rather than
converging on a shared physical mechanism, the literature remains
bifurcated between small-area electrostatic optimization and the
thermo-chemo-mechanical realities of operational modules.

\subsection{10. Reporting Practice Over Six
Months}\label{reporting-practice-over-six-months}

A fundamental disagreement in inverted perovskite solar cells (PSCs)
concerns how self-assembled monolayer (SAM) hole-selective contacts
degrade during operation: while one investigation attributes operational
failure to out-of-plane molecular desorption from metal oxides that
rigid planar conjugation suppresses
{[}\href{https://doi.org/10.1016/j.scib.2026.02.050}{163}{]}, another
concludes that degradation proceeds via field-induced lateral
aggregation into isolated clusters
{[}\href{https://doi.org/10.1038/s41467-026-72097-1}{165}{]}. Resolving
whether molecular detachment or lateral dewetting governs contact
failure requires operando bias-dependent sum-frequency generation
spectroscopy paired with in situ conductive atomic force microscopy
under continuous maximum power point tracking (MPPT) at operational
temperatures. Reconciling such conflicting degradation mechanisms across
the literature is hindered by the scarcity of standardized testing
protocols: when stability assessments rely on disparate stress regimes,
distinguishing true interfacial stabilization from divergent testing
conditions becomes impossible.

Across the six months, pooled reporting rates over 2238 indexed works
reveal substantial deficits in primary performance and reliability
metrics. Only 52.1\% (1167 of 2238) of abstracts stated any power
conversion efficiency (PCE) value. Independent verification was rare:
only 8.5\% (191 of 2238) reported an independently certified efficiency,
leaving most performance claims susceptible to unverified spectral
mismatch or contact overestimation. Stabilised performance tracking
appeared in only 11.0\% (246 of 2238) of abstracts, while reporting of
current-voltage hysteresis was confined to 2.7\% (60 of 2238), allowing
capacitive transients and ion accumulation to obscure steady-state power
output. Device geometry was similarly omitted, with only 8.4\% (188 of
2238) stating the active area. Standardized operational aging was almost
entirely absent, with only 1.7\% (38 of 2238) identifying an
International Summit on Organic Photovoltaic Stability (ISOS) protocol
label. For a reader attempting cross-device comparisons, these low
detection rates mean that nominally equivalent efficiencies often
represent fundamentally incomparable physical measurements. These
detection rates reflect automated parsing of published abstracts rather
than an exhaustive audit of experimental methodologies; omitting a
testing protocol or geometric parameter in an abstract does not
establish that the underlying investigation failed to perform the
measurement.

Examining monthly ranges demonstrates that reporting practices did not
improve across the half year. Stating an efficiency sat flat within a
narrow monthly range of 48.1\% to 55.3\%. Reporting of certified
efficiencies oscillated between 5.9\% and 11.6\%, while stabilised
efficiency remained confined to a tight monthly range of 9.2\% to
13.5\%. Active area disclosure sat flat between 5.8\% and 10.3\%,
hysteresis identification stayed low within 1.5\% to 4.3\%, and ISOS
labeling remained marginal between 0.6\% and 2.7\%. Although the monthly
share of closely read papers expanded steadily from January (12.6\%),
February (15.0\%), and March (16.7\%) through April (18.8\%), May
(17.7\%), and June (19.2\%), this growing literature volume produced no
upward drift in reporting completeness. Practice did not evolve;
instead, reporting metrics occupied fixed corridors throughout the
entire six months.

Exemplary reporting occurs where full geometric, certification, and
operational parameters converge. Large-area module studies provide
verified baselines by simultaneously stating aperture dimensions,
certified efficiencies, and standardized degradation conditions,
including an active area of \qty{100}{\centi\meter\squared} with a
certified PCE of 25.55\% under ISOS-L-1 tracking
{[}\href{https://doi.org/10.1126/sciadv.aec3238}{13}{]} and
\qty{100}{\centi\meter\squared} with a certified PCE of 23.55\% under
ISOS-D-1 storage
{[}\href{https://doi.org/10.1038/s41467-026-69685-6}{14}{]},
establishing reproducible benchmarks for scalable interconnects
{[}\href{https://doi.org/10.1126/sciadv.aec3238}{13},\href{https://doi.org/10.1038/s41467-026-69685-6}{14}{]}.
Standardized operational testing is also illustrated by small cells
identifying an ISOS-L-2 protocol and certified PCE of 21.54\% alongside
a time to 80\% of initial performance (T80) of \qty{1210}{\hour}
{[}\href{https://doi.org/10.1038/s41467-026-72099-z}{207}{]}, continuous
operational stress under ISOS-L-3
{[}\href{https://doi.org/10.1021/acsenergylett.5c03999}{174}{]}, or
ISOS-L2 testing with a certified 26.78\% PCE
{[}\href{https://doi.org/10.1002/adma.202520432}{173}{]}, providing
unambiguous tracking of photothermal decay
{[}\href{https://doi.org/10.1002/adma.202520432}{173},\href{https://doi.org/10.1021/acsenergylett.5c03999}{174},\href{https://doi.org/10.1038/s41467-026-72099-z}{207}{]}.
In sharp contrast, absent reporting obscures physical mechanisms.
Extreme claims, such as an uncertified champion PCE of 28.49\%
{[}\href{https://doi.org/10.1007/s40820-025-02006-6}{208}{]}, omit
certification, active area, and standardized aging protocols from the
abstract. High champion efficiencies are frequently reported on
sub-\qty{0.1}{\centi\meter\squared} cells, such as
\qty{0.072}{\centi\meter\squared} at 26.2\% PCE
{[}\href{https://doi.org/10.1038/s41467-025-68207-0}{209}{]} or
\qty{0.058}{\centi\meter\squared} at 26.87\% PCE
{[}\href{https://doi.org/10.1002/aenm.70925}{169}{]}, where lateral
current collection and edge encapsulation distort planar transport
physics
{[}\href{https://doi.org/10.1002/aenm.70925}{169},\href{https://doi.org/10.1038/s41467-025-68207-0}{209}{]}.
At the opposite extreme, scaling studies describe an aperture of
\qty{900}{\centi\meter\squared} without providing an independently
certified efficiency
{[}\href{https://doi.org/10.1002/adfm.202532205}{18}{]}. Uncertified
claims evaluated on tiny cells without standardized protocols impede
physical comparison and leave operational longevity unanchored.

\subsection{11. Research Gaps and
Outlook}\label{research-gaps-and-outlook}

A fundamental disagreement in multi-junction investigations concerns
whether non-radiative voltage deficits originate from defective charge
extraction at buried contacts or from trap-assisted recombination within
bulk wide-bandgap (WBG) absorbers. Anchoring self-assembled monolayers
(SAMs) directly to a reactive plasma deposition (RPD) tin oxide (SnOx)
recombination layer to suppress contact barriers and halide accumulation
delivers a certified power conversion efficiency (PCE) of 33.6\%
(champion 33.2\% on \qty{1}{\centi\meter\squared}) and certified 31.0\%
on a \qty{207.9}{\centi\meter\squared} minimodule
{[}\href{https://doi.org/10.1126/science.aef5355}{1}{]}. In direct
contradiction, competing reports conclude that carrier losses are
dominated by photoinduced halide segregation and deep trap formation
within the bulk WBG film, asserting that polymeric polycarbazole
phosphonic acid (Poly-3Ph-4PACz) suppresses bulk recombination to yield
a certified two-terminal (2T) tandem PCE of 32.95\% (champion 26.13\%,
stabilised 27.12\% in single junctions; 22.22\% on a
\qty{57.6}{\centi\meter\squared} flexible module)
{[}\href{https://doi.org/10.1038/s41467-026-73276-w}{2}{]}. Standard
current-voltage scans cannot resolve this conflict because contact
recombination velocity and bulk quasi-Fermi level splitting collapse
produce identical open-circuit voltage (Voc) deficits. Decoupling
contact resistance from bulk non-radiative recombination requires
operando, spatially resolved voltage-dependent photoluminescence (PL)
imaging via the local transport and recombination study (LTRS) method
paired with bias-dependent transient absorption spectroscopy under
continuous illumination to isolate interfacial from bulk carrier losses
{[}\href{https://doi.org/10.1002/smll.202513958}{138}{]}.

A second unresolved contradiction governs operational degradation at
perovskite interfaces, where studies diverge on whether hole-selective
contact decay proceeds via vertical molecular detachment or lateral
dewetting, and whether bulky passivators relieve lattice strain or
accelerate fatigue. One investigation attributes contact failure to
out-of-plane molecular desorption that planar conjugated backbones
suppress {[}\href{https://doi.org/10.1016/j.scib.2026.02.050}{163}{]},
whereas another concludes that degradation is driven by
electric-field-induced lateral migration of self-assembled monolayer
(SAM) molecules into isolated clusters
{[}\href{https://doi.org/10.1038/s41467-026-72097-1}{165}{]}.
Compounding this ambiguity, bulky organic passivators designed to raise
phase transition barriers, such as 3-PMPCl reaching a certified 27.6\%
PCE in formamidinium lead triiodide (\ce{FAPbI3})
{[}\href{https://doi.org/10.1126/science.aeb9953}{159}{]}, induce local
lattice expansion under thermal stress at heterointerfaces, driving
interfacial fatigue even while unencapsulated devices retain 87\%
initial PCE after \qty{4092}{\hour} of maximum power point tracking
(MPPT) under International Summit on Organic Photovoltaic Stability
(ISOS) protocol ISOS-L-1 with a time to 80\% of initial performance
(T80) of \qty{3566}{\hour}
{[}\href{https://doi.org/10.1126/sciadv.aef6596}{162}{]}. Deciding
whether contact failure originates from covalent bond cleavage, lateral
dewetting, or thermal expansion mismatch requires operando
bias-dependent sum-frequency generation spectroscopy coupled with in
situ high-temperature grazing-incidence wide-angle X-ray scattering
under continuous electrical load.

A third research gap concerns conflicting kinetic pathways for
stabilizing the photoactive black alpha phase of \ce{FAPbI3} against
polymorphic transition into the photoinactive yellow delta phase.
Contact-triggered cationic interaction asserts that phase stability
occurs through non-destructive physical contact between separately
crystallized two-dimensional (2D) and three-dimensional (3D) perovskites
without interdiffusion, yielding a certified 25.61\% PCE (champion
26.25\%) {[}\href{https://doi.org/10.1038/s41560-026-02027-4}{98}{]}. In
direct contrast, recrystallization using formamidinium tetrafluoroborate
demonstrates that forming stable mixed-dimensional networks requires
active formamidinium/methylammonium cation exchange, while 2D
perovskitoid interlayers disrupt delta-phase nucleation to achieve an
uncertified champion 25.61\% PCE
{[}\href{https://doi.org/10.1038/s41467-026-73739-0}{100}{]}.
Concurrently, low-dimensional passivators reach certified efficiencies
of 26.2\% (champion 26.49\%)
{[}\href{https://doi.org/10.1038/s41467-026-72790-1}{19}{]}, but whether
these capping layers introduce resistive transport bottlenecks remains
unproven. Resolving whether phase stabilization proceeds by passive
interfacial templating or dynamic chemical exchange requires in situ
time-resolved grazing-incidence X-ray scattering paired with operando
nanoscale secondary ion mass spectrometry during thermal annealing and
light soaking.

A fourth mechanistic gap lies in opposing kinetic models of
precursor-solvent coordination during scalable meniscus coating. One
framework contends that prolonged solvent-solute interaction during
blade coating stabilizes detrimental solvent-adduct intermediates,
demanding rapid adduct suppression to achieve 26.5\% PCE on
\qty{0.0665}{\centi\meter\squared}
{[}\href{https://doi.org/10.1038/s41467-026-68439-8}{186}{]}.
Conversely, competing investigations conclude that strong coordination
by N-methyl-2-pyrrolidone or dimethyl sulfoxide is indispensable to
retard premature nucleation and stabilize intermediate complexes,
improving Voc from 1.08 to \qty{1.09}{\volt} and fill factor (FF) from
75\% to 80.5\% {[}\href{https://doi.org/10.1063/5.0297564}{187}{]}, or
achieving certified 26.22\% PCE on \qty{0.06}{\centi\meter\squared}
through polypropylene glycol microcrystal regulation
{[}\href{https://doi.org/10.1002/adma.202522508}{194}{]}. However, these
coordination models remain demonstrated only on
sub-\qty{0.1}{\centi\meter\squared} apertures without verifying that
intermediate equilibria persist across large areas under industrial
drying and mechanical shear
{[}\href{https://doi.org/10.20517/ss.2025.128}{195}{]}. Closing this gap
requires in situ optical reflectance and photoluminescence spectroscopy
coupled with synchrotron X-ray diffraction during high-speed slot-die
coating across modules exceeding \qty{100}{\centi\meter\squared}.

Underpinning these physical gaps is a persistent empirical disconnect
between certified peak efficiency, device area, and operational
durability verification. Across the half year, the best certified PCE
reached 34.85\% in a 2T tandem, yet only 114 of 872 closely read papers
reported any certified value. The best certified value moved from 29.2\%
in January to 33.6\% in June, a change of 4.4 percentage points over the
half year, while uncertified claims outpaced certified benchmarks,
reaching 41.6\% in April and 44.36\% in May. Durability verification
remains similarly sparse: only 39 papers reported a T80 lifetime, with
the longest reaching \qty{4100}{\hour}, and only 39 papers named an ISOS
protocol. Although the largest reported aperture reached
\qty{900}{\centi\meter\squared}, champion efficiencies remain confined
to sub-\qty{0.1}{\centi\meter\squared} cells that mask sheet resistance,
interconnect losses, and drying gradients. Crucially, this half year's
evidence could not support a rigorous comparison between buried contact
passivators and bulk wide-bandgap stabilizers because reported devices
diverged across uncalibrated active areas and unstandardized aging
regimes. Instead of unencapsulated room-temperature scans on
micro-cells, the field must measure encapsulated devices under
continuous operational load using standardized ISOS-L-2 or ISOS-L-3
protocols at 65 °C, reporting certified steady-state power output on
apertures of at least \qty{1}{\centi\meter\squared} alongside spatially
resolved luminescence maps of local contact resistance.

\subsection{Acknowledgements}\label{acknowledgements}

The authors would like to acknowledge the financial support from Xiamen
University Malaysia through the Xiamen University Malaysia Research Fund
(XMUMRF/2022-C10/IENG/0046), and from National Natural Science
Foundation of China (Ref. No.: 52503390). The authors would also like to
thank Anthropic for support through the Anthropic AI for Science
program.

\subsection{Declaration of Competing
Interest}\label{declaration-of-competing-interest}

The authors declare no conflict of interest.

\subsection{Data Availability
Statement}\label{data-availability-statement}

The corpus table for the month, the per-paper extraction records with
their verbatim source quotations, the statistics file underlying every
number in the text, and the gate report are provided as Supplementary
Information. Abstract text and full texts are not redistributed.

\subsection{AI Usage Declaration}\label{ai-usage-declaration}

This manuscript was prepared as part of our Research Project for making
a reliable Autonomous Researcher Agent. The agent executed the monthly
review pipeline without human intervention: it harvested the indexed
literature for the month through the OpenAlex, Semantic Scholar,
Crossref and arXiv interfaces, applied a deterministic topical filter,
ranked and selected the papers read closely, extracted every
quantitative claim into a structured record binding each number to a
verbatim quotation from that paper's own abstract, computed the
statistics and figures from those records, and drafted this text from
the resulting evidence briefs. Literature harvesting, selection,
statistics, figures and gate checking are deterministic scripts rather
than model output. Structured extraction used qwen3.8-flash through the
opencode interface; the prose was written by gemini-3.8-flash-high
through the agy interface; independent review used Opus 5 at high
reasoning effort. No figure is a model-generated image: every panel is
plotted directly from the extracted records. Every quantitative
statement in the main text resolves to the machine-generated statistics
file shipped as Supplementary Information, and every citation resolves
to a record carrying its source quotation. The author reviewed and
edited the content, verified the selection, and takes full
responsibility for the content of the published article.

\subsection{References}\label{references}

\begin{enumerate}
\def\labelenumi{\arabic{enumi}.}
\tightlist
\item
  Shi, W. et al.~Indium-free perovskite/silicon tandem solar cells with
  tin oxide recombination layer and electrodes. \emph{Science} (2026).
  \url{https://doi.org/10.1126/science.aef5355}
\item
  Tang, Q. et al.~Buried-interface homogenization by asymmetric
  polymeric self-assembled layers powers efficient, durable flexible
  perovskite photovoltaics. \emph{Nature Communications} (2026).
  \url{https://doi.org/10.1038/s41467-026-73276-w}
\item
  Zhang, X. et al.~Minimizing the VOC deficit by regulating facets of
  wide-bandgap perovskite for monolithic perovskite/silicon tandem solar
  cells. \emph{eScience} (2026).
  \url{https://doi.org/10.1016/j.esci.2026.100563}
\item
  Parida, B. et al.~Solution and Vacuum‐Based Scalable Deposition
  Methods for Perovskite/Si Tandem Solar Cells. \emph{Advanced Optical
  Materials} (2026). \url{https://doi.org/10.1002/adom.202503341}
\item
  Zheng, J. et al.~In Situ Engineered Gradient Ruddlesden--Popper 2D/3D
  Heterostructures at the Buried Interface Enable Current-Matched
  Perovskite/Silicon Tandems. \emph{ACS Energy Letters} (2026).
  \url{https://doi.org/10.1021/acsenergylett.5c04127}
\item
  Chen, Y. et al.~Stabilizing Efficient Perovskite/Silicon Tandem Solar
  Cells With Amorphous and Asymmetric Hole‐Selective Contacts.
  \emph{Advanced Materials} (2026).
  \url{https://doi.org/10.1002/adma.73106}
\item
  Du, H. et al.~Carbon-incorporated polysilicon interconnection layer
  enables robust self-assembled monolayer anchoring for
  perovskite/TOPCon tandem solar cells. \emph{Nature Communications}
  (2026). \url{https://doi.org/10.1038/s41467-026-72794-x}
\item
  Luo, M. et al.~Uniform and ordered self-assembled hole-selective
  layers driven by π-π interactions for efficient perovskite-silicon
  tandem solar cells. \emph{Nature Communications} (2026).
  \url{https://doi.org/10.1038/s41467-026-72160-x}
\item
  Jiang, X. et al.~Regulating intermolecular interaction of passivators
  for controllable surface energetics of photovoltaic perovskites.
  \emph{Science Advances} (2026).
  \url{https://doi.org/10.1126/sciadv.aec4431}
\item
  Yang, J. et al.~Modulating Buried Interfaces via Co‐Assembled
  Monolayers With Siloxane‐Mediated Intermolecular Interactions for
  Perovskite/Silicon Tandem Solar Cells. \emph{Advanced Energy
  Materials} (2026). \url{https://doi.org/10.1002/aenm.71107}
\item
  Yan, T. et al.~Sterically-extended asymmetric conjugated
  hole-selective layer for perovskite/silicon tandem solar cells.
  \emph{Nature Communications} (2026).
  \url{https://doi.org/10.1038/s41467-026-73656-2}
\item
  He, Q. et al.~Molecular Pinning Effect in Composite Self‐Assembled
  Monolayers Enabling Efficient Perovskite/TOPCon Tandem Solar Cells.
  \emph{Angewandte Chemie International Edition} (2026).
  \url{https://doi.org/10.1002/anie.5162616}
\item
  Zhong, Y. et al.~Homogenizing out-of-plane strain distribution for
  high-performance flexible perovskite photovoltaics. \emph{Science
  Advances} (2026). \url{https://doi.org/10.1126/sciadv.aec3238}
\item
  Xie, Y. et al.~Regulating perovskite crystallization kinetics at laser
  scribe lines for efficient and stable perovskite modules. \emph{Nature
  Communications} (2026).
  \url{https://doi.org/10.1038/s41467-026-69685-6}
\item
  Huang, S. et al.~Molecularly Engineered Self‐Healing Scaffold With
  Customized Dynamic Bonds Enable Stable and Scalable Flexible
  Perovskite Solar Cells and Modules. \emph{Advanced Materials} (2026).
  \url{https://doi.org/10.1002/adma.73855}
\item
  Wang, D. et al.~Constructing an Ultra‐Efficient Carrier Transport
  Pathway Through Dual‐Targeted Anchoring for High‐Performance
  Perovskite Solar Cells. \emph{Advanced Energy Materials} (2026).
  \url{https://doi.org/10.1002/aenm.71111}
\item
  Liu, J. et al.~In‐situ Dual‐stage Reactions Enhance Operational
  Stability of High‐Efficiency Perovskite Modules. \emph{Solar RRL}
  (2026). \url{https://doi.org/10.1002/solr.70353}
\item
  Hwang, S.J. et al.~Auxeticity‐by‐Assembly: Interlocking Modular
  Auxetic Metamaterials with Selectively Activatable AgNW--Graphene
  Oxide‐EGaIn Composite Interconnects for Scalable Freeform Photovoltaic
  Modules. \emph{Advanced Functional Materials} (2026).
  \url{https://doi.org/10.1002/adfm.202532205}
\item
  Liu, Y. et al.~Liquid-derived, solvent-free vapor-mediated dimensional
  reconstruction yields a record fill factor in inverted perovskite
  solar cells. \emph{Nature Communications} (2026).
  \url{https://doi.org/10.1038/s41467-026-72790-1}
\item
  Wang, Y. et al.~In Situ Formation of Guanidinium Derivative and 2D/3D
  Heterostructure for 26.49\% Efficient and Stable Inverted Perovskite
  Solar Cells. \emph{Angewandte Chemie International Edition} (2026).
  \url{https://doi.org/10.1002/anie.8890696}
\item
  Wu, C. et al.~Defect Passivation with Organic Co‐Crystal Layers for
  Efficient and Stable Inverted Perovskite Solar Cells: A
  Non‐2D‐Perovskites Strategy for Surface Treatments. \emph{Advanced
  Materials} (2026). \url{https://doi.org/10.1002/adma.202519138}
\item
  Yang, T. et al.~Stereoelectronic manipulation of ligands for
  perovskite solar cells. \emph{Nature} (2026).
  \url{https://doi.org/10.1038/s41586-026-10626-0}
\item
  Yan, Y. et al.~Imide‐Based Molecular Passivation Enables Efficient
  Wide‐Bandgap Perovskite and Perovskite/Silicon Tandem Solar Cells.
  \emph{Advanced Energy Materials} (2026).
  \url{https://doi.org/10.1002/aenm.70902}
\item
  Hou, W. et al.~Lewis Acid--Base Coordination‐Driven Interface
  Passivation Using 3‐Hydroxyflavone for Inverted Perovskite Solar
  Cells. \emph{Small} (2026). \url{https://doi.org/10.1002/smll.73673}
\item
  Su, W. et al.~Antisolvent Engineering via Fluorinated Phenylalkylamine
  Additives for Synergistic Optimization on Crystallization Kinetics and
  Carrier Transport of Perovskite Solar Cells. \emph{ACS Applied Energy
  Materials} (2026). \url{https://doi.org/10.1021/acsaem.6c00562}
\item
  Dai, Y. et al.~Redox‐Driven Spatiotemporal Passivation for Efficient
  and Durable Perovskite Solar Cells. \emph{Advanced Functional
  Materials} (2026). \url{https://doi.org/10.1002/adfm.76623}
\item
  Wei, F. et al.~Efficient Rigid and Flexible Perovskite Photovoltaics
  with Crystallization and Morphology Evolution via Chlorine-Containing
  Triazine Molecular Passivation. \emph{ACS Materials Letters} (2026).
  \url{https://doi.org/10.1021/acsmaterialslett.5c01678}
\item
  Pan, W. et al.~Defect-Anchored Dipole Molecules Induce Surface
  Polarization Facilitating High-Performance Inverted Perovskite Solar
  Cells. \emph{Nano-Micro Letters} (2026).
  \url{https://doi.org/10.1007/s40820-026-02150-7}
\item
  Chu, W. et al.~Halogen‐Functional Molecules for Synergistic
  Multi‐Defect Healing in High‐Performance Perovskite Solar Cells.
  \emph{Small} (2026). \url{https://doi.org/10.1002/smll.73436}
\item
  Zhang, K. et al.~Synergistic Passivation Strategy for Efficient and
  Stable Perovskite Solar Cells. \emph{Advanced Functional Materials}
  (2026). \url{https://doi.org/10.1002/adfm.202532110}
\item
  Yang, R. et al.~Quantitative Deposition Enables Dual Passivation
  Synergy for Efficient Inverted Perovskite Solar Cells. \emph{Small}
  (2026). \url{https://doi.org/10.1002/smll.73886}
\item
  Wang, Y. et al.~Synergistic Bulk and Interfacial Defect Control Boosts
  Voltage and Efficiency in Perovskite/Silicon Tandem Solar Cells.
  \emph{ACS Applied Materials \& Interfaces} (2026).
  \url{https://doi.org/10.1021/acsami.6c01711}
\item
  Wang, R. et al.~Chiral Molecular Configuration‐Mediated
  Crystallization of Perovskites for High‐Performance Solar Cells.
  \emph{Advanced Energy Materials} (2026).
  \url{https://doi.org/10.1002/aenm.71222}
\item
  Qi, Y. et al.~Eliminated Interfacial Side Reactions in Perovskite
  Solar Cells by Sterically Protected Ammonium Passivation. \emph{ACS
  Applied Materials \& Interfaces} (2026).
  \url{https://doi.org/10.1021/acsami.6c04956}
\item
  Zhang, L. et al.~Unlocking Synergistic Passivation in Perovskite Solar
  Cells via Molecular Design of Electron‐Donation and Steric Hindrance.
  \emph{Advanced Functional Materials} (2026).
  \url{https://doi.org/10.1002/adfm.75743}
\item
  Pi, J. et al.~Temporal Decoupling of Processing and Passivation: A
  Latent Passivator Activated by In Situ Hydrolysis for Efficient
  Perovskite Solar Cells. \emph{Advanced Functional Materials} (2026).
  \url{https://doi.org/10.1002/adfm.75370}
\item
  Sun, P. et al.~From Literature to Lab: Closed-Loop Advancement of
  Perovskite Solar Cells via Domain Knowledge Guided LLM. \emph{arXiv
  (Cornell University)} (2026).
  \url{https://doi.org/10.48550/arxiv.2602.04914}
\item
  Y, Z. et al.~Molecularly Confined Domains Enable Halide‐Stable
  Wide‐Bandgap Perovskites. \emph{Advanced Materials} (2026).
  \url{https://doi.org/10.1002/adma.73713}
\item
  Liu, R.X. et al.~Migration-Resistant Phenylethylguanidinium Salts
  Stabilizing Perovskite Solar Cells. \emph{ACS Energy Letters} (2026).
  \url{https://doi.org/10.1021/acsenergylett.6c00737}
\item
  Zhao, W. et al.~Polymerized Surface Passivation for Stable and
  Efficient Inverted Perovskite Solar Cells. \emph{Advanced Materials}
  (2026). \url{https://doi.org/10.1002/adma.202522825}
\item
  Lan, Y. et al.~Synergistic Molecular Locking Through Sodium‐Integrated
  Cross‐Linkable Scaffold Enables Durable Perovskite Solar Cells and
  Modules. \emph{Angewandte Chemie International Edition} (2026).
  \url{https://doi.org/10.1002/anie.202525815}
\item
  Shi, N. et al.~Multifunctional Fluorinated Self-Polymerizable Additive
  Improves the Performance of Perovskite Photovoltaics. \emph{ACS Energy
  Letters} (2026). \url{https://doi.org/10.1021/acsenergylett.5c03901}
\item
  Zhao, R. et al.~In Situ Interfacial Polymerization Enabling a
  Dual‐Anchor Surface Binding Interlayer for Efficient and Stable
  Inverted Perovskite Solar Cells. \emph{Angewandte Chemie International
  Edition} (2026). \url{https://doi.org/10.1002/anie.7579788}
\item
  Jiang, C. et al.~Suppressing Interlayer Ion Migration With a Tailored
  Passivator for Highly Air‐Stable Perovskite Solar Cells.
  \emph{Advanced Functional Materials} (2026).
  \url{https://doi.org/10.1002/adfm.75412}
\item
  Xiong, J. et al.~Managing Grain Boundary Energy by Dipole Strategy in
  Sn‐Pb Perovskite for All‐Perovskite Tandems. \emph{Advanced Materials}
  (2026). \url{https://doi.org/10.1002/adma.72755}
\item
  Xu, X. et al.~Mechanochemically Reinforced Dual‐Dynamic Covalent
  Seeding Enables High‐Performance and Operationally Stable Perovskite
  Solar Cells. \emph{Advanced Materials} (2026).
  \url{https://doi.org/10.1002/adma.72823}
\item
  Zhang, Z. et al.~Photoswitchable isomers to improve grain boundary
  resilience and perovskite solar cells stability under light cycling.
  \emph{Nature Energy} (2026).
  \url{https://doi.org/10.1038/s41560-026-01993-z}
\item
  Hu, B. et al.~Viologen Enabled Dynamic Interface Passivation for
  Efficient and Stable Inverted Perovskite Solar Cells. \emph{Advanced
  Functional Materials} (2026).
  \url{https://doi.org/10.1002/adfm.202530317}
\item
  Shen, H. et al.~Homogenizing interfacial assembly via indole-mediated
  binary monolayers for perovskite solar cells. \emph{Nature
  Communications} (2026).
  \url{https://doi.org/10.1038/s41467-026-71301-6}
\item
  Deng, Z. et al.~A Co‐Self‐Assembled Monolayer With Enhanced Coverage
  and Defect Passivation for Efficient Inverted Perovskite Solar Cells.
  \emph{Energy \& environment materials} (2026).
  \url{https://doi.org/10.1002/eem2.70443}
\item
  Shao, C. et al.~Rationally Designed Self‐Assembled Monolayer for
  Dual‐Site Passivation Enables Efficient and Stable Perovskite Solar
  Cells. \emph{Advanced Materials} (2026).
  \url{https://doi.org/10.1002/adma.202523249}
\item
  Fu, X. et al.~Defect‐Passivating and Dense Indolocarbazole‐Based
  Self‐Assembled Monolayers for Efficient Inverted Perovskite Solar
  Cells With over 26.1\% Efficiency. \emph{Small} (2026).
  \url{https://doi.org/10.1002/smll.202512942}
\item
  Wei, K. et al.~Multidentate Molecular Anchoring for Enhanced
  Interfacial Stability and Reliable Perovskite Solar Cells.
  \emph{Advanced Energy Materials} (2026).
  \url{https://doi.org/10.1002/aenm.70806}
\item
  Ge, C. et al.~Stabilizing Sputtered NiO x via In Situ Dissociative
  Adsorption Passivation for Efficient Perovskite Solar Cells.
  \emph{Advanced Materials} (2026).
  \url{https://doi.org/10.1002/adma.73382}
\item
  Li, B. et al.~Simple‐Synthesis Isomer‐Free Multi‐Adduct Fullerenes as
  Electron Transport Materials Enable 26.66\% Efficiency of Perovskite
  Solar Cells. \emph{Angewandte Chemie International Edition} (2026).
  \url{https://doi.org/10.1002/anie.4500167}
\item
  Jiang, Y. et al.~Fused-Ring Electron Acceptors as a Versatile Additive
  Platform for Efficient Perovskite Photovoltaics. \emph{Journal of the
  American Chemical Society} (2026).
  \url{https://doi.org/10.1021/jacs.5c16976}
\item
  Feng, Z. et al.~Managing Crystallization in Tin‐Lead Perovskite Solar
  Cells via Differential Chelating Coordination. \emph{Energy \&
  environment materials} (2026).
  \url{https://doi.org/10.1002/eem2.70453}
\item
  You, J. et al.~Synergistic Management of Crystallization and Defects
  in Tin‐Lead Perovskites towards Efficient All‐Perovskite Tandem Solar
  Cells. \emph{Advanced Materials} (2026).
  \url{https://doi.org/10.1002/adma.73038}
\item
  Wang, Y. et al.~Halide‐Cation Interactions Enable Controlled
  Crystallization and Defect Minimization in High‐Performance
  All‐Perovskite Tandems. \emph{Advanced Materials} (2026).
  \url{https://doi.org/10.1002/adma.202522668}
\item
  Yousuf, H. et al.~Annealing‐Induced Interface Passivation in n‐TOPCon
  Silicon Solar Cells for High‐Efficiency 4T Perovskite/Si Tandems.
  \emph{Energy Technology} (2026).
  \url{https://doi.org/10.1002/ente.70501}
\item
  Mishra, K.P. et al.~Engineering Solar Cell Interfaces With
  Transition‐Metal Dichalcogenides: Insights From SCAPS‐1D Simulations.
  \emph{Energy Technology} (2026).
  \url{https://doi.org/10.1002/ente.70553}
\item
  Liu, L. et al.~Electric Field Distribution Engineering at the Buried
  Interface for High‐Performance Inverted Perovskite Solar Cells.
  \emph{Advanced Functional Materials} (2026).
  \url{https://doi.org/10.1002/adfm.76291}
\item
  Wang, J. et al.~Asymmetric Self‐Assembled Molecule With High Dipole
  Moment for Efficient Wide‐Bandgap Perovskite Solar Cells and Tandems.
  \emph{Advanced Materials} (2026).
  \url{https://doi.org/10.1002/adma.202523338}
\item
  Zhang, X. et al.~Tuning Terminal Groups of Hole‐Selective Materials to
  Suppress Nonradiative Recombination in Inverted Perovskite Solar
  Cells. \emph{Angewandte Chemie International Edition} (2026).
  \url{https://doi.org/10.1002/anie.8473744}
\item
  Gong, S. et al.~Dynamic Feedback of Interfacial Electric Fields
  Governs Charge Extraction in Perovskite Solar Cells. \emph{Advanced
  Functional Materials} (2026). \url{https://doi.org/10.1002/adfm.75295}
\item
  Yao, S. et al.~Sulfur‐Modulated Fully‐Conjugated Self‐Assembled
  Monolayers for Synergistic Dual‐Interface Optimization in Inverted
  Perovskite Solar Cells. \emph{Angewandte Chemie International Edition}
  (2026). \url{https://doi.org/10.1002/anie.5319815}
\item
  Wang, Y. et al.~Molecular Polarization‐Driven Synergistic Interface
  Engineering for High‐Performance Perovskite Solar Cells.
  \emph{Advanced Materials} (2026).
  \url{https://doi.org/10.1002/adma.73306}
\item
  Zhou, Z. et al.~Molecular Bridge Regulation of Buried Interface in
  Perovskite Solar Cells. \emph{Advanced Materials} (2026).
  \url{https://doi.org/10.1002/adma.202519267}
\item
  Chen, X. et al.~Photo‐Tautomerization‐Driven Energy Transfer at the
  Hole‐Transport Interface Stabilizes Efficient Inverted Perovskite
  Solar Cells. \emph{Angewandte Chemie International Edition} (2026).
  \url{https://doi.org/10.1002/anie.3596437}
\item
  Wang, W. et al.~Atomic Contact Enables Highly Oriented and Covered
  Self-Assembled Monolayers for Composite Hole-Selective Structure
  Achieving Efficient Perovskite Solar Cells. \emph{ACS Applied
  Materials \& Interfaces} (2026).
  \url{https://doi.org/10.1021/acsami.6c03372}
\item
  Wei, X. et al.~Interface Engineering With Annealed MXene and SAM
  Molecules for High‐Performance Inverted Perovskite Solar Cells.
  \emph{Advanced Energy Materials} (2026).
  \url{https://doi.org/10.1002/aenm.70966}
\item
  Liu, Z. et al.~Kinetic Enhancement of Self-Assembled Molecular
  Deposition via Catalytic Additive for Inverted Perovskite Solar Cells.
  \emph{ACS Energy Letters} (2026).
  \url{https://doi.org/10.1021/acsenergylett.6c01257}
\item
  Zhang, J. et al.~Area‐Selective Atomic Layer Deposition of AlO x at
  the Buried Interface for High‐Performance Perovskite Solar Cells.
  \emph{Angewandte Chemie International Edition} (2026).
  \url{https://doi.org/10.1002/anie.202516537}
\item
  Wang, J. et al.~Redox‐Mediated Stabilization of the Hole Transport
  Layer and Buried Interface Toward Stable Perovskite Solar Cells.
  \emph{Angewandte Chemie International Edition} (2026).
  \url{https://doi.org/10.1002/anie.4012708}
\item
  Pang, L.Y. et al.~Dual Functional Molecular Engineering Utilizes
  Synergistic Effects to Enhance Hole Transport in Inverted Perovskite
  Solar Cells. \emph{Advanced Energy Materials} (2026).
  \url{https://doi.org/10.1002/aenm.71005}
\item
  Guo, S. et al.~Vertically Oriented Bisphosphonate Self‐Assembled
  Monolayers for Efficient and Thermocycling‐Stable Inverted Perovskite
  Solar Cells. \emph{Advanced Energy Materials} (2026).
  \url{https://doi.org/10.1002/aenm.70814}
\item
  Geng, Y. et al.~Synthesis of a Butterfly‐Shaped π‐Extended
  Dibenzazepine and Its Application in Interfacial Engineering for
  High‐Performance Perovskite Solar Cells. \emph{Small} (2026).
  \url{https://doi.org/10.1002/smll.74278}
\item
  Yu, X. et al.~Self‐Assembly Molecular Ordering for Strengthened
  Interface and Efficient Perovskite/Silicon Tandem Solar Cells.
  \emph{Small} (2026). \url{https://doi.org/10.1002/smll.74319}
\item
  Wu, W. et al.~Synergistic Density Enhancement and Energetic Alignment
  via Tailored Postassembly of Self-Assembled Molecules for Perovskite
  Solar Cells. \emph{The Journal of Physical Chemistry Letters} (2026).
  \url{https://doi.org/10.1021/acs.jpclett.6c01042}
\item
  Hou, M. et al.~Mitigating Photothermal‐Induced Burn‐In in Perovskite
  Solar Cells via a Multidentate Phosphonic Acid Polymer Network.
  \emph{Advanced Energy Materials} (2026).
  \url{https://doi.org/10.1002/aenm.71159}
\item
  Zhang, Y. et al.~Thermo‐Mechanically Robust Inverted Perovskite Solar
  Cells Toward Extreme Environments. \emph{Advanced Functional
  Materials} (2026). \url{https://doi.org/10.1002/adfm.75697}
\item
  Huang, K. et al.~Enhancing Crystallization and Optimizing the Buried
  Interface via Hydrophobic Trifluoromethyl Molecules for Efficient
  Sn--Pb Perovskite and All‐Perovskite Tandem Solar Cells.
  \emph{Advanced Energy Materials} (2026).
  \url{https://doi.org/10.1002/aenm.70910}
\item
  Xiao, Y. et al.~N‐Heterocycle‐Activated π‐Cooperative Coordination for
  Enhancing Structural Stability of Perovskite Solar Cells.
  \emph{Advanced Materials} (2026).
  \url{https://doi.org/10.1002/adma.73828}
\item
  Wang, F. et al.~Enhancing Heterogeneous Nucleation on Buried Interface
  for Efficient Antisolvent‐Free Inverted Flexible Perovskite
  Photovoltaics. \emph{Advanced Materials} (2026).
  \url{https://doi.org/10.1002/adma.73358}
\item
  Xiong, Z. et al.~Precursor-State Preorganization Narrows
  Crystallization Pathways in Formamidinium-Rich Perovskite Solar Cells.
  \emph{ACS Energy Letters} (2026).
  \url{https://doi.org/10.1021/acsenergylett.6c01291}
\item
  Jiang, W. et al.~Manipulating Buried Interface via Aromatic Amino Acid
  Derivatives for Wide‐Bandgap Perovskite Solar Cells and Tandem
  Devices. \emph{Small} (2026). \url{https://doi.org/10.1002/smll.73595}
\item
  Yuan, S. et al.~Unified Interface--Bulk Enhancements in Wide‐Bandgap
  Perovskites Enabled by Bifunctional Cyanoacetic Acid for Textured
  Silicon/Perovskite Tandem Solar Cells. \emph{Advanced Materials}
  (2026). \url{https://doi.org/10.1002/adma.72738}
\item
  Yang, Y. et al.~Multifunctional Conjugated Ligands Synergistically
  Optimize the Interface and Regulate Crystallization for
  High‐Performance 2D/3D Perovskite Solar Cells. \emph{Advanced
  Materials} (2026). \url{https://doi.org/10.1002/adma.202520251}
\item
  Wang, H. et al.~Multifunctional Thiazolidine Derivatives Synergistic
  Regulation of Charge Transport Interface for Inverted Perovskite Solar
  Cells. \emph{ACS Energy Letters} (2026).
  \url{https://doi.org/10.1021/acsenergylett.5c03328}
\item
  Ye, S. et al.~The Bimolecular Co‐Modification Strategy for Interface
  Management Enables High‐Performance Inverted Perovskite Solar Cells
  Without Pre‐Deposited HTL. \emph{Advanced Functional Materials}
  (2026). \url{https://doi.org/10.1002/adfm.76198}
\item
  Peng, S. et al.~In Situ Polymerizable Lead-Chelating Crown Ether for
  Lead Leakage Inhibition in Perovskite Solar Cells. \emph{ACS Energy
  Letters} (2026). \url{https://doi.org/10.1021/acsenergylett.6c00703}
\item
  Yang, Y. et al.~Hafnium oxide interface stabilization for efficient,
  photothermally stable perovskite solar cells. \emph{Science} (2026).
  \url{https://doi.org/10.1126/science.aea3339}
\item
  Wang, K. et al.~Robust self-assembled monolayer enables ultraviolet
  stable perovskite photovoltaics. \emph{Nature Communications} (2026).
  \url{https://doi.org/10.1038/s41467-026-73426-0}
\item
  Gao, J. et al.~Synergistic Dual‐Passivation of Grain Boundaries and
  Buried Interface for High‐Efficiency and Stable Perovskite Solar
  Cells. \emph{Small} (2026). \url{https://doi.org/10.1002/smll.73891}
\item
  Xu, B. et al.~Interface‐Triggered Bulk Electrical Coupling in SnO 2
  via Mesoscopic Chemical Reconfiguration for Scalable, High‐Efficiency
  Perovskite Photovoltaics. \emph{Advanced Materials} (2026).
  \url{https://doi.org/10.1002/adma.73286}
\item
  Taneja, Y. et al.~Computational and Experimental Investigations of
  Mg‐Doped SnO 2 Electron Transport Layer for Indoor Perovskite Solar
  Cells. \emph{Small} (2026).
  \url{https://doi.org/10.1002/smll.202512808}
\item
  Liu, Q. et al.~Electronic Blockade of Shunting Pathways via Dual
  Insulator Contacts for High-Efficiency Wide-Bandgap Perovskite Indoor
  Photovoltaics. \emph{Nano-Micro Letters} (2026).
  \url{https://doi.org/10.1007/s40820-026-02225-5}
\item
  Lee, S.H. et al.~Contact-triggered molecular interactions enable
  structural refinement of perovskite layers in solar cells.
  \emph{Nature Energy} (2026).
  \url{https://doi.org/10.1038/s41560-026-02027-4}
\item
  Shi, J. et al.~Constructing a Soft‐Hard Interpenetrating Network in
  Flexible Quasi‐2D Perovskite Solar Cells Through FABF 4 ‐Driven Phase
  Redistribution. \emph{Advanced Functional Materials} (2026).
  \url{https://doi.org/10.1002/adfm.75673}
\item
  Xing, C. et al.~Mechanistic insight into α-to-δ phase transition and
  stabilizing α-phase FAPbI3 via regulating δ-phase orientation.
  \emph{Nature Communications} (2026).
  \url{https://doi.org/10.1038/s41467-026-73739-0}
\item
  Dong, H. et al.~Deprotonation suppressing via competitive proton
  transfer control for efficient perovskite solar cells. \emph{Nature
  Communications} (2026).
  \url{https://doi.org/10.1038/s41467-026-73620-0}
\item
  Wang, Z. et al.~Citric Acid‐Inhibited Phase Separation Assists In Situ
  Growth of Ternary Cationic Perovskite Single Crystal Thin Films Toward
  Efficient Solar Cells. \emph{Advanced Energy Materials} (2026).
  \url{https://doi.org/10.1002/aenm.70939}
\item
  Zheng, L. et al.~Enabling Highly Efficient and Stable Perovskite
  Photovoltaics via A Multidentate Molecular Anchor Additive.
  \emph{Nano-Micro Letters} (2026).
  \url{https://doi.org/10.1007/s40820-026-02098-8}
\item
  Fu, Z. et al.~Cation Volatility Governs the Functionality of
  Thiocyanate Additives in FAPbI3 Perovskite Solar Cells. \emph{ACS
  Applied Energy Materials} (2026).
  \url{https://doi.org/10.1021/acsaem.6c01329}
\item
  Fu, Y. et al.~Complete DMA + /Cs + Exchange and Rapid Phase Conversion
  via FAAC-Mediated Intermediate Phase Engineering for Efficient and
  Stable CsPbI 3 Perovskite Solar Cells. \emph{The Journal of Physical
  Chemistry Letters} (2026).
  \url{https://doi.org/10.1021/acs.jpclett.6c01395}
\item
  Wang, Z. et al.~Facile Molecular Modification for Efficient and Stable
  Carbon-Based CsPbI 3 Perovskite Solar Cells. \emph{ACS Applied Energy
  Materials} (2026). \url{https://doi.org/10.1021/acsaem.5c04098}
\item
  Huang, W. et al.~Solvent‐Additive Synergy Optimizes the Performance of
  Methylamine‐Based Printable Mesoscopic Perovskite Solar Cells.
  \emph{Solar RRL} (2026). \url{https://doi.org/10.1002/solr.70363}
\item
  Wang, X. et al.~LEAP: A closed-loop framework for perovskite precursor
  additive discovery. \emph{arXiv (Cornell University)} (2026).
  \url{https://doi.org/10.48550/arxiv.2605.20242}
\item
  Rao, Z. et al.~Multiple Synergistic Mechanisms of Urea in
  Vacuum‐Assisted Prepared Triple‐Cation Perovskite Solar Cells.
  \emph{Solar RRL} (2026). \url{https://doi.org/10.1002/solr.202500906}
\item
  Xiong, G. et al.~Anion‐Engineered Reinforcement of Hydrogen Bonding
  toward Homogeneous and Stable Perovskite Solar Cells. \emph{Advanced
  Functional Materials} (2026). \url{https://doi.org/10.1002/adfm.76197}
\item
  Li, S.Y. et al.~Regulation of Crystallization Kinetics via
  Zwitterionic Additive for Efficient and Operationally Stable
  Formamidinium Perovskite Solar Cells. \emph{Advanced Energy Materials}
  (2026). \url{https://doi.org/10.1002/aenm.71105}
\item
  Li, B.Z. et al.~All Air‐Processed n‐i‐p Perovskite Solar Cells Exceed
  26\% Efficiency Enabled by Thiacetazone Additive. \emph{Advanced
  Materials} (2026). \url{https://doi.org/10.1002/adma.202523630}
\item
  Wu, J. et al.~Approaching 26\% Efficiency in Inverted FAPbI 3
  Perovskite Solar Cells Enabled by Tailored Fluoropyridine Derivative
  Additives. \emph{Advanced Functional Materials} (2026).
  \url{https://doi.org/10.1002/adfm.202526679}
\item
  Li, W. et al.~Reversible crosslinking strategy for dynamic strain
  regulation in inverted perovskite solar cells. \emph{Nature
  Communications} (2026).
  \url{https://doi.org/10.1038/s41467-026-70697-5}
\item
  Liu, F. et al.~Misplaced-dipole engineered repairable fluoropolymer
  elastomer for flexible perovskite solar cell with excellent
  thermal-mechanical cycling resistance. \emph{Nature Communications}
  (2026). \url{https://doi.org/10.1038/s41467-026-74285-5}
\item
  Yeo, J.H. et al.~Dry Vacuum Processing of Perovskite Films: Insights
  Into Growth and Precision Control Enabling Batch‐Stable
  High‐Efficiency Solar Cells. \emph{Carbon Energy} (2026).
  \url{https://doi.org/10.1002/cey2.70236}
\item
  Qin, R. et al.~Eutectic Urea Phosphate Suppresses Phase Segregation in
  Wide‐Bandgap Photovoltaic Perovskites. \emph{Angewandte Chemie
  International Edition} (2026).
  \url{https://doi.org/10.1002/anie.2535666}
\item
  Gu, T. et al.~Enhancing Halogen Phase Synchronous Crystallization
  Enables High‐Performance Wide‐Bandgap Perovskite Photovoltaics.
  \emph{Small} (2026). \url{https://doi.org/10.1002/smll.202513111}
\item
  Xu, Z. et al.~Synergistic A-Site Codoping Enables Lattice
  Stabilization and Strain Balance in Wide-Bandgap Perovskites.
  \emph{ACS Energy Letters} (2026).
  \url{https://doi.org/10.1021/acsenergylett.6c00328}
\item
  Heydarian, M. et al.~Chloride‐Substituted Perovskite Enables Thick
  Middle Cell for Perovskite/Perovskite/Silicon Triple‐Junction Solar
  Cells. \emph{Solar RRL} (2026).
  \url{https://doi.org/10.1002/solr.70410}
\item
  Su, H. et al.~Targeted Functional Group‐Modulated Benzoic Additive for
  High‐Performance 2‐Terminal Perovskite/Silicon Tandems. \emph{Advanced
  Functional Materials} (2026). \url{https://doi.org/10.1002/adfm.75640}
\item
  Zheng, H. et al.~Multisite‐Functionalized Thiazole Spacer Enables
  Synergistic Grain‐Boundary Healing and Energy‐Level Alignment for
  High‐Performance 2D/3D Perovskite Solar Cells. \emph{Advanced
  Functional Materials} (2026).
  \url{https://doi.org/10.1002/adfm.202524302}
\item
  Liu, P. et al.~Synergistic Interface Engineering via a Dual-A-Site 1D
  Perovskite Bridge for Efficient All-Perovskite Tandems. \emph{ACS
  Energy Letters} (2026).
  \url{https://doi.org/10.1021/acsenergylett.5c04243}
\item
  Utomo, D.S. et al.~One-Step Formation of 2D/3D Perovskite
  Heterojunction via Ligand Intercalation and Facet Engineering for
  Efficient Perovskite Solar Cells. \emph{Nano-Micro Letters} (2026).
  \url{https://doi.org/10.1007/s40820-025-02058-8}
\item
  Y, L. et al.~A Diammonium‐Based Non‐Dion‐Jacobson Phase 2D Perovskite
  With High Durability for Efficient and Stable 2D/3D Perovskite Solar
  Modules. \emph{Angewandte Chemie International Edition} (2026).
  \url{https://doi.org/10.1002/anie.2974414}
\item
  Liu, M. et al.~Constructing Layered 1D/3D Heterostructure for
  Efficient and Stable Single‐Crystal Perovskite Photovoltaics.
  \emph{Advanced Functional Materials} (2026).
  \url{https://doi.org/10.1002/adfm.75386}
\item
  Tan, H. et al.~Strain Regulation through In Situ 1D/3D Heterojunction
  Construction with an Optimal Alkyl Chain Length for Efficient and
  Stable Perovskite Solar Cells. \emph{ACS Applied Energy Materials}
  (2026). \url{https://doi.org/10.1021/acsaem.6c01363}
\item
  Li, G. et al.~Grain Boundary Encapsulation for Efficient Perovskite
  Solar Cells via PbI 2 ‐Derived Low‐Dimensional Structure.
  \emph{Advanced Functional Materials} (2026).
  \url{https://doi.org/10.1002/adfm.202531852}
\item
  Tang, X. et al.~Methoxyl‐Substitution of Phenylethylammonium
  Strengthened 2D/3D Heterogeneous Structures for Inverted Perovskite
  Solar Cells with Enhanced Efficiency and Stability. \emph{Small}
  (2026). \url{https://doi.org/10.1002/smll.73351}
\item
  Feng, W. et al.~Breaking Electronic Insulation of Monocyclic Aromatic
  Spacers via Hydrazide‐Induced Orbital Coupling in Ruddlesden‐Popper
  Perovskites. \emph{Angewandte Chemie International Edition} (2026).
  \url{https://doi.org/10.1002/anie.4293157}
\item
  Jiang, B. et al.~Highly‐Oriented 3D Perovskite Films Synthesized by
  the Assisted Growth of Large‐Area 2D Perovskite Nanosheets With n = 1.
  \emph{Small} (2026). \url{https://doi.org/10.1002/smll.74076}
\item
  Lang, Y. et al.~Two-Dimensional Perovskite Quantum Dot Building Blocks
  toward Efficient 2D/3D Perovskite Heterostructures in Photovoltaics.
  \emph{ACS Applied Materials \& Interfaces} (2026).
  \url{https://doi.org/10.1021/acsami.5c24915}
\item
  Xu, B. et al.~Impact of Functional Group Configuration in Isomeric
  Additives on Device Performance of Quasi-2D Perovskite Solar Cells.
  \emph{The Journal of Physical Chemistry Letters} (2026).
  \url{https://doi.org/10.1021/acs.jpclett.6c01531}
\item
  Ning, F. et al.~Additive‐Mediated (111)‐Facet‐Oriented Perovskite
  Films via Two‐Step Blade‐Coating for Efficient and Stable Inverted
  Solar Cells. \emph{Solar RRL} (2026).
  \url{https://doi.org/10.1002/solr.70385}
\item
  Gayen, K. et al.~Enhancing the Performance of Perovskite Solar Cells
  through Interfacial Modification by Fluorinated 2D Perovskite.
  \emph{ACS Applied Energy Materials} (2026).
  \url{https://doi.org/10.1021/acsaem.6c00285}
\item
  Huang, S. et al.~Vacuum Co‐Evaporation for High Reproducible and
  Environmentally Stable Hybrid Perovskite Thin Film Solar Cells by
  Cation Doping. \emph{Energy Technology} (2026).
  \url{https://doi.org/10.1002/ente.202502217}
\item
  Wang, Q. et al.~Engineering a Conformal and Lossless Interconnect via
  Ultrathin ITO for High‐Efficiency Perovskite/Silicon Tandem Solar
  Cells. \emph{Advanced Functional Materials} (2026).
  \url{https://doi.org/10.1002/adfm.75311}
\item
  Fischer, O. et al.~Quantifying the Spatial Distribution of Series
  Resistance in Monolithic Perovskite/Silicon Tandem Solar Cells Using
  Voltage‐Dependent Photoluminescence Imaging. \emph{Small} (2026).
  \url{https://doi.org/10.1002/smll.202513958}
\item
  Huang, G. et al.~Impact of Ultrathin Electron-Selective Contacts on
  Perovskite--Silicon Tandem Solar Cells. \emph{ACS Energy Letters}
  (2026). \url{https://doi.org/10.1021/acsenergylett.5c04082}
\item
  Shi, B. et al.~Antimony oxide buffer layer for single- and
  double-junction perovskite-based solar cells. \emph{Nature
  Communications} (2026).
  \url{https://doi.org/10.1038/s41467-026-70848-8}
\item
  Xue, X. et al.~Dopant-Free Ultra-Thin Spiro-OMeTAD Enables Near
  30\%-Efficient n--i--p Perovskite/Silicon Tandem Solar Cells.
  \emph{Nano-Micro Letters} (2026).
  \url{https://doi.org/10.1007/s40820-025-01962-3}
\item
  Jiang, W. et al.~Stable Self‐Assembled Monolayers Anchoring and
  Enhanced Antireflection in Four‐Terminal Perovskite/Silicon Tandem
  Solar Cells via a VO x ‐NiO x Composite Contact. \emph{Advanced Energy
  Materials} (2026). \url{https://doi.org/10.1002/aenm.71036}
\item
  Huang, N. et al.~Pathway-priming of intermediate phases in FACs-based
  wide-bandgap perovskites for Si tandem solar cells. \emph{Nature
  Communications} (2026).
  \url{https://doi.org/10.1038/s41467-026-72094-4}
\item
  Mondal, S. et al.~Impact of Br and Cs Substitution on Optoelectronic
  Properties in FAPbI3 for Perovskite/TOPCon Tandem Solar Cells using
  First-Principles and Drift-Diffusion Simulation Study. \emph{ChemRxiv}
  (2026). \url{https://doi.org/10.26434/chemrxiv.15001694/v1}
\item
  Tian, R. et al.~Chemical hardness engineering synchronizes
  crystallization in perovskite tandems. \emph{Nature Nanotechnology}
  (2026). \url{https://doi.org/10.1038/s41565-026-02165-6}
\item
  He, Z. et al.~Volatile Fluorination Additive for Operationally Stable
  Perovskite/TOPCon Tandem Solar Cells. \emph{Advanced Functional
  Materials} (2026). \url{https://doi.org/10.1002/adfm.202526009}
\item
  Huang, Q. et al.~Coordination Engineering Enables Efficient 1.77 eV
  Wide‐Bandgap Perovskites for All‐Perovskite Tandem Solar Cells.
  \emph{Solar RRL} (2026). \url{https://doi.org/10.1002/solr.70387}
\item
  Kim, S. et al.~Selective Ion‐Blocking Strategy for Stabilizing
  Perovskite Tandem Solar Cells. \emph{Advanced Energy Materials}
  (2026). \url{https://doi.org/10.1002/aenm.202506634}
\item
  Mao, Y. et al.~Halogen Bonding Networks Regulate Bromide
  Crystallization and Stabilize Phase Homogeneity in Efficient
  All‐Perovskite Tandem Solar Cells. \emph{Advanced Functional
  Materials} (2026). \url{https://doi.org/10.1002/adfm.74659}
\item
  Han, L. et al.~Crystallization Modulation and Halide Segregation
  Suppression in Wide‐Bandgap Perovskite for Efficient All‐Perovskite
  Tandem Solar Cells. \emph{Small} (2026).
  \url{https://doi.org/10.1002/smll.202510652}
\item
  Chai, W. et al.~Surface Reconstruction Through Synergetic
  PEACl‐Modification and Solvent Etching for High‐Efficiency and Stable
  Perovskite/Silicon Tandem Solar Cells. \emph{Advanced Energy
  Materials} (2026). \url{https://doi.org/10.1002/aenm.202506619}
\item
  Wang, Q. et al.~A Universal Strategy to Stabilize Crystal Lattice and
  Composition of Wide‐Bandgap Perovskites for Efficient Tandem Solar
  Cells. \emph{Advanced Materials} (2026).
  \url{https://doi.org/10.1002/adma.72986}
\item
  Lu, Z. et al.~Targeted Enhancement of Passivator Coordination Activity
  Enables 2.00 V V oc in Efficient Perovskite/Silicon Tandem Solar
  Cells. \emph{Advanced Energy Materials} (2026).
  \url{https://doi.org/10.1002/aenm.71077}
\item
  Guo, Y. et al.~Stoichiometry-Regulated Crystallization and Halide
  Homogenization in Wide-Bandgap Perovskites for Efficient Solar Cells
  and Tandems. \emph{ACS Nano} (2026).
  \url{https://doi.org/10.1021/acsnano.6c04740}
\item
  Kuba, A.G. et al.~Bandgap Tunable Two-Step Vapor-Deposited Perovskite
  Absorbers for Perovskite-Silicon Tandem Solar Cells. \emph{ACS Energy
  Letters} (2026). \url{https://doi.org/10.1021/acsenergylett.6c00156}
\item
  Mardegan, L. et al.~Sequential Evaporation for Scalable Hybrid
  Processing of Perovskite/Silicon Tandem Solar Cells. \emph{ACS Applied
  Materials \& Interfaces} (2026).
  \url{https://doi.org/10.1021/acsami.6c07582}
\item
  Zhou, X. et al.~Synchronous Interfacial and Crystallographic
  Engineering of Wide-Bandgap Perovskites for Efficient
  Perovskite/Silicon Tandem Solar Cells. \emph{ACS Energy Letters}
  (2026). \url{https://doi.org/10.1021/acsenergylett.6c00166}
\item
  Li, W. et al.~Redox Cascade‐Driven Structural Degradation Under
  Reverse Bias in All‐Perovskite Tandem Solar Cells. \emph{Advanced
  Materials} (2026). \url{https://doi.org/10.1002/adma.202522942}
\item
  Miao, T. et al.~An entropy-regulating molecular lock stabilizes
  formamidinium lead halide perovskite. \emph{Science} (2026).
  \url{https://doi.org/10.1126/science.aeb9953}
\item
  Huang, W. et al.~Interference Crystallization Rebalances Facet
  Competition for Efficient and Stable Perovskite Solar Cells.
  \emph{Advanced Materials} (2026).
  \url{https://doi.org/10.1002/adma.73539}
\item
  Yang, Y. et al.~Molecular Buffering Regulates Lattice Strain for
  Fatigue‐Resistant Perovskite Photovoltaics Under Cryogenic Thermal
  Cycling. \emph{Advanced Materials} (2026).
  \url{https://doi.org/10.1002/adma.202519339}
\item
  Ren, Z. et al.~Reactive reconstruction and embedded passivation of
  heterointerfaces for intrinsically stable perovskite photovoltaics.
  \emph{Science Advances} (2026).
  \url{https://doi.org/10.1126/sciadv.aef6596}
\item
  Li, D. et al.~Strongly spatial-confined self-assembled monolayers for
  high-performance perovskite photovoltaics. \emph{中国科学通报：英文版}
  (2026). \url{https://doi.org/10.1016/j.scib.2026.02.050}
\item
  Cen, T. et al.~Strengthening Substrate Anchoring of Polymeric
  Self‐Assembled Monolayers for Efficient and Stable Inverted Perovskite
  Solar Cells. \emph{Advanced Materials} (2026).
  \url{https://doi.org/10.1002/adma.72622}
\item
  Zhuang, R. et al.~Co-deposited inverted perovskite photovoltaics
  towards 27\% efficiency via vertical redistribution of
  self-assembled-molecules and in-situ crosslinking. \emph{Nature
  Communications} (2026).
  \url{https://doi.org/10.1038/s41467-026-72097-1}
\item
  Yang, G. et al.~Molecular confinement for enhanced interfacial contact
  and performance in inverted perovskite solar cells. \emph{Nature
  Communications} (2026).
  \url{https://doi.org/10.1038/s41467-026-72115-2}
\item
  Gu, Z. et al.~2D Perovskite Engineering Enables Robust Self‐Assembled
  Monolayers for High‐Performance Perovskite Solar Cells. \emph{Advanced
  Materials} (2026). \url{https://doi.org/10.1002/adma.73561}
\item
  Wang, X. et al.~Hydroxyl Interfacial Engineering for Self‐Assemble
  Monolayers Anchoring on NiO x Enables Efficient and Stable Perovskite
  Solar Cells. \emph{Advanced Materials} (2026).
  \url{https://doi.org/10.1002/adma.73656}
\item
  Wu, X. et al.~Amino Acid Hydrochloride‐Mediated Self‐Assembled
  Monolayers for High‐Performance Inverted Perovskite Solar Cells.
  \emph{Advanced Energy Materials} (2026).
  \url{https://doi.org/10.1002/aenm.70925}
\item
  Yang, Z. et al.~Spatial Molecular Engineering of Hole Semiconductors
  Enables Record Efficiency and Durability in Inverted Perovskite Solar
  Cells. \emph{Angewandte Chemie International Edition} (2026).
  \url{https://doi.org/10.1002/anie.202523665}
\item
  Xiong, Q. et al.~Quantifying Deep‐Level Defects‐Dominated Degradation
  for Commercially Viable Perovskite Solar Cells. \emph{Advanced
  Materials} (2026). \url{https://doi.org/10.1002/adma.73679}
\item
  Chang, Q. et al.~Thermotropic liquid-assisted interface management
  enables efficient and stable perovskite solar cells and modules.
  \emph{Nature Communications} (2026).
  \url{https://doi.org/10.1038/s41467-025-68231-0}
\item
  Huang, Z. et al.~Positional Methylation Isomer Additives Regulate the
  Degradation Pathway of Perovskites. \emph{Advanced Materials} (2026).
  \url{https://doi.org/10.1002/adma.202520432}
\item
  Yang, Q. et al.~Dual-Site Competitive Coordination Eliminates Residual
  DMSO for Thermally and Moisture-Stable Perovskite Solar Cells.
  \emph{ACS Energy Letters} (2026).
  \url{https://doi.org/10.1021/acsenergylett.5c03999}
\item
  Yang, Q. et al.~Regulating the Crystallization‐Defect‐Strain Improves
  the Stability and Scalability of FACs Perovskite Photovoltaics.
  \emph{Advanced Functional Materials} (2026).
  \url{https://doi.org/10.1002/adfm.76756}
\item
  Liu, X. et al.~Conjugated Amine Ligand‐Induced Robust‐Conductive
  Interlayer Enables Efficient and Durable n‐i‐p Perovskite Solar Cells.
  \emph{Advanced Materials} (2026).
  \url{https://doi.org/10.1002/adma.73161}
\item
  Yu, Q. et al.~Molecular Suppression of PCBM Dimerization in Inverted
  Perovskite Solar Cells. \emph{Solar RRL} (2026).
  \url{https://doi.org/10.1002/solr.70405}
\item
  Li, Q. et al.~Multifunctional Buffer Layer for Bolstering the
  Stability and Photovoltaic Performance of Perovskite Solar Cells.
  \emph{Advanced Materials} (2026).
  \url{https://doi.org/10.1002/adma.73529}
\item
  Hou, E. et al.~Kilogram-scale one-pot synthesis of multicomponent
  fullerene composites for efficient inverted perovskite solar cells.
  \emph{Nature Communications} (2026).
  \url{https://doi.org/10.1038/s41467-026-70022-0}
\item
  Li, Z. et al.~Interfacial Energetics Reconstruction via Bridging
  Engineering for Efficient Inverted Perovskite Solar Cells and Modules.
  \emph{Angewandte Chemie International Edition} (2026).
  \url{https://doi.org/10.1002/anie.9286497}
\item
  Kim, Y. et al.~Controlling Doping Preference of Spiro‐OMeTAD for
  Efficient and Stable Perovskite Solar Cells. \emph{Advanced Energy
  Materials} (2026). \url{https://doi.org/10.1002/aenm.202506760}
\item
  Liu, H. et al.~Ionic Anchoring of Dopants via a Multifunctional
  Phenothiazine Regulator Enables Perovskite Solar Cells With Enhanced
  Operational Longevity. \emph{Advanced Functional Materials} (2026).
  \url{https://doi.org/10.1002/adfm.76257}
\item
  Zhang, H. et al.~Lithium‐Based Deep Eutectic in Spiro‐OMeTAD Enable
  Efficient and Stable Perovskite Solar Cells. \emph{Angewandte Chemie
  International Edition} (2026).
  \url{https://doi.org/10.1002/anie.9551164}
\item
  Wang, Z. et al.~Chemical Polishing‐Assisted Surface Passivation and
  Repair for Robust Perovskite Photovoltaics. \emph{Advanced Functional
  Materials} (2026). \url{https://doi.org/10.1002/adfm.76110}
\item
  Wang, C.F. et al.~Unlocking the Full Potential of Polydentate
  Self‐Assembled Molecules in Perovskite Solar Cells: Co‐Deposition
  Versus Pre‐Deposition. \emph{Advanced Functional Materials} (2026).
  \url{https://doi.org/10.1002/adfm.74893}
\item
  Jin, L. et al.~Suppressing solvent adducts via coordination
  competition enables scalable perovskite photovoltaics. \emph{Nature
  Communications} (2026).
  \url{https://doi.org/10.1038/s41467-026-68439-8}
\item
  Bai, X. et al.~DMF-free co-solvent doctor-bladed FA0.6MA0.4PbI3
  perovskite solar cells and modules under ambient conditions.
  \emph{Applied Physics Letters} (2026).
  \url{https://doi.org/10.1063/5.0297564}
\item
  Pusapati, R. et al.~Improved Solar Cell Performance of 1.68 eV
  Perovskite Absorber Through Scaffold and Solvent Engineering for
  Scalable Two‐Step Hybrid Deposition. \emph{Solar RRL} (2026).
  \url{https://doi.org/10.1002/solr.70406}
\item
  Singh, J. et al.~Sustainable ternary solvent strategy for
  high-efficiency slot-die coated perovskite solar cells.
  \emph{Materials Reports Energy} (2026).
  \url{https://doi.org/10.1016/j.matre.2026.100406}
\item
  Miyake, S. et al.~Giant‐Domain Perovskite Films via
  Reduced‐Temperature Meniscus‐Guided Coating: Fluid‐Dynamic Control of
  Supersaturation. \emph{Solar RRL} (2026).
  \url{https://doi.org/10.1002/solr.202600003}
\item
  Li, M. et al.~A Novel Transformer Architecture for Scalable Perovskite
  Thin-Film Detection. \emph{Micromachines} (2026).
  \url{https://doi.org/10.3390/mi17030314}
\item
  Peng, L. et al.~High-Performance Perovskite Solar Cells and Modules
  via Slot-Die Coating: Engineering Advances, Photoluminescence
  Insights, and Future Perspectives. \emph{Nano-Micro Letters} (2026).
  \url{https://doi.org/10.1007/s40820-026-02134-7}
\item
  Li, Z. et al.~Nanoscale Colloids Engineering for Minimizing Efficiency
  Loss in Scalable Perovskite Solar Cells. \emph{Advanced Materials}
  (2026). \url{https://doi.org/10.1002/adma.202520776}
\item
  Shen, F. et al.~Scalable Fabrication of High‐Efficiency Thin‐Film
  Perovskite Solar Cells in Air via Polymer‐Mediated Synthesis of
  α‐FAPbI 3 Microcrystals. \emph{Advanced Materials} (2026).
  \url{https://doi.org/10.1002/adma.202522508}
\item
  Zong, H. et al.~Chemical-mechanical co-design for scalable flexible
  perovskite manufacturing. \emph{Soft Science} (2026).
  \url{https://doi.org/10.20517/ss.2025.128}
\item
  Zhang, Y. et al.~Toward scalable solar cells: optimizing
  atmospheric-pressure spatial atomic layer deposition SnOX for enhanced
  perovskite crystallization and performance. \emph{Solar Energy
  Materials and Solar Cells} (2026).
  \url{https://doi.org/10.1016/j.solmat.2026.114233}
\item
  Xu, H.W. et al.~Efficient Perovskite Solar Modules Enabled by
  Synergistic Interparticle and Interfacial Electrostatic Interaction
  Regulation. \emph{ACS Energy Letters} (2026).
  \url{https://doi.org/10.1021/acsenergylett.6c00476}
\item
  Hongsith, K. et al.~Scalable fabrication of SnO2 quantum dot
  electron-transport layers through multicycle infrared annealing for
  perovskite solar cells. \emph{Results in Engineering} (2026).
  \url{https://doi.org/10.1016/j.rineng.2026.111690}
\item
  Du, H. et al.~Interface-Directed Self-Assembled Monolayers for
  High-Efficiency and Scalable Inverted Wide-Bandgap Perovskite Solar
  Cells. \emph{ACS Applied Materials \& Interfaces} (2026).
  \url{https://doi.org/10.1021/acsami.6c03670}
\item
  Li, Z. et al.~Low-hygroscopic solvents enable ambient blade coating of
  efficient perovskite solar cells. \emph{Nature Communications} (2026).
  \url{https://doi.org/10.1038/s41467-026-72581-8}
\item
  Vanin, F. et al.~Stable comprehensive dopants for efficient and
  scalable perovskite solar cells. \emph{Joule} (2026).
  \url{https://doi.org/10.1016/j.joule.2025.102310}
\item
  Gong, C. et al.~Coffee‐Ring‐Engineered Interfacial Layer Enables
  Fluid‐Locking for Scalable Perovskite Photovoltaics. \emph{Advanced
  Materials} (2026). \url{https://doi.org/10.1002/adma.202518582}
\item
  Salim, K.M.M. et al.~Defect Suppression and Efficiency Enhancement of
  Slot-Die-Coated Formamidinium Lead Iodide Solar Devices via Ambient
  Air Annealing. \emph{ACS Applied Materials \& Interfaces} (2026).
  \url{https://doi.org/10.1021/acsami.5c22032}
\item
  Fang, Y. et al.~Scalable solution soaking quenching technique unlocks
  efficient and durable wide bandgap perovskite solar modules.
  \emph{Nature Communications} (2026).
  \url{https://doi.org/10.1038/s41467-026-69264-9}
\item
  Kumar, N. et al.~Paradigm shift in perovskite photovoltaics
  scalability: Spray-engineered 1D/3D perovskite interfaces unlock
  efficiency and moisture stability. \emph{Materials Today
  Sustainability} (2026).
  \url{https://doi.org/10.1016/j.mtsust.2026.101347}
\item
  Huang, S. et al.~A Defect‐Engineered Vacuum Evaporation Strategy for
  High‐Efficiency Indoor Perovskite Mini Solar Modules. \emph{Small
  Methods} (2026). \url{https://doi.org/10.1002/smtd.70654}
\item
  Zhang, H. et al.~Self-assembled 1D/3D heterojunction enables
  all-inorganic perovskite 4-terminal tandem solar cells with 21.54\%
  certified efficiency. \emph{Nature Communications} (2026).
  \url{https://doi.org/10.1038/s41467-026-72099-z}
\item
  Meng, R. et al.~Nature-Inspired Redox Shuttle with Regenerable
  Antioxidant for Efficient All-Perovskite Tandem Solar Cells.
  \emph{Nano-Micro Letters} (2026).
  \url{https://doi.org/10.1007/s40820-025-02006-6}
\item
  Du, K. et al.~Anti-aggregation self-assembled monolayers enable
  high-performance and scalable perovskite solar cells. \emph{Nature
  Communications} (2026).
  \url{https://doi.org/10.1038/s41467-025-68207-0}
\end{enumerate}

\end{document}
